% file: manual.tex
% Copyright 2008-2013 V. Bos, T. van Deursen, P. Kordy, and S. Mauw
% This file is part of the MSC Macro Package.
%
\documentclass[12pt,a4paper]{article}
\def\pgfautoxrefs{1} 
\usepackage[version=latest]{pgf}
\usepackage{xparse}
% \usepackage{a4wide}
% \usepackage{url}
\usepackage{moreverb}
\usepackage{multicol}
\usepackage{msc}
\usepackage{showexpl}
\usepackage{xkeyval,calc,listings,tikz,fp}
\usepackage{hyperref}
\usepackage{ifluatex}
\usepackage[a4paper,left=2.25cm,right=2.25cm,top=2.5cm,bottom=2.5cm,nohead]{geometry}
\usepackage{amsmath,amssymb}
\usepackage{xxcolor}
\usepackage{pifont}
\usepackage{makeidx}
\usepackage{compsci}
\usepackage{xstring}
% \usepackage{pgfmanual}
\newcommand{\PK}[1]{%
{\color{red}{\bf[PK]} #1}}

% we allow a ragged right
\setlength{\rightskip}{0pt plus 0.05\linewidth minus 0pt}

% \newcommand{\cmd}[1]{\texttt{\bslash #1}}

\newcommand{\acro}[1]{{#1}}

\newcommand{\MSC}{\acro{MSC}}
\newcommand{\HMSC}{\acro{HMSC}}
\newcommand{\MSCdoc}{\MSC{}doc}
\newcommand{\mscpack}{\MSC{} macro package}
% \newcommand{\env}[1]{\texttt{#1}}
\newcommand{\cmdarg}[1]{\{\emph{#1}\}}
\newcommand{\coordarg}[1]{\emph{#1}}
\newcommand{\coordargs}[2]{(\coordarg{#1},\coordarg{#2})}
\newcommand{\lnsvalue}[3]{large/normal/small value #1/#2/#3}

\newenvironment{defs}{%
  \begin{list}{}%
              {\setlength{\labelwidth}{0pt}%
               \setlength{\labelsep}{1em}%
               \setlength{\leftmargin}{1em}%
               \setlength{\parsep}{1ex}%
               \setlength{\listparindent}{0pt}%
               \setlength{\rightmargin}{0pt}%
               \renewcommand{\makelabel}[1]{##1}%
               \raggedright%
              }%
  }{%
  \end{list}}
\newcommand{\tikzname}{Ti\emph{k}Z}
\newcommand{\beamer}{\textsc{beamer}}
\newcommand{\pdf}{\textsc{pdf}}
\newcommand{\eps}{\texttt{eps}}
\newcommand{\pgfname}{\textsc{pgf}}
% \newcommand{\prosper}{\textsc{prosper}}
% \newcommand{\seminar}{\textsc{seminar}}
% \newcommand{\texpower}{\textsc{texpower}}
% \newcommand{\foils}{\textsc{foils}}

% \providecommand\href[2]{\texttt{#1}}
% \providecommand\hypertarget[2]{\texttt{#1}}
% \providecommand\hyperlink[2]{\texttt{#1}}


\colorlet{examplefill}{yellow!80!black}
\definecolor{graphicbackground}{rgb}{1,1,1}
\definecolor{codebackground}{rgb}{0.90,0.90,1}

\newenvironment{pgfmanualentry}{\par\vspace*{-8pt}\list{}{\leftmargin=2em\itemindent-\leftmargin\def\makelabel##1{\hss##1}}}{\endlist\par\vspace*{-8pt}}
\newcommand\pgfmanualentryheadline[1]{\itemsep=0pt\parskip=0pt\item\strut{#1}\par\topsep=0pt}
\newcommand\pgfmanualbody{\parskip3pt\vspace*{-3pt}}



\newenvironment{pgflayout}[1]{
  \begin{pgfmanualentry}
    \pgfmanualentryheadline{\texttt{\string\pgfpagesuselayout\char`\{\declare{#1}\char`\}}\oarg{options}}
    \index{#1@\protect\texttt{#1} layout}%
    \index{Page layouts!#1@\protect\texttt{#1}}%
    \pgfmanualpdflabel{#1}{}%
    \pgfmanualbody
}
{
  \end{pgfmanualentry}
}


\newenvironment{command}[1]{
  \begin{pgfmanualentry}
    \extractcommand#1\@@
    \pgfmanualbody
}
{
  \end{pgfmanualentry}
}

\newenvironment{commandlist}[1]{%
  \begin{pgfmanualentry}
  \foreach \xx in {#1} {%
    \expandafter\extractcommand\xx\@@
  }%
  \pgfmanualbody
}{%
  \end{pgfmanualentry}
}%

% \begin{internallist}[register]{\pgf@xa}
% \end{internallist}
%
% \begin{internallist}[register]{\pgf@xa,\pgf@xb}
% \end{internallist}
\newenvironment{internallist}[2][register]{%
  \begin{pgfmanualentry}
  \foreach \xx in {#2} {%
    \expandafter\extractinternalcommand\expandafter{\xx}{#1}%
  }%
  \pgfmanualbody
}{%
  \end{pgfmanualentry}
}%
\def\extractinternalcommand#1#2{%
  \pgfmanualentryheadline{Internal #2 \declare{\texttt{\string#1}}}%
  \removeats{#1}%
  \index{Internals!\strippedat @\protect\myprintocmmand{\strippedat}}%
  \index{\strippedat @\protect\myprintocmmand{\strippedat}}%
  \pgfmanualpdflabel{\textbackslash\strippedat}{}%
}

%% MW: START MATH MACROS
\def\mvar#1{{\ifmmode\textrm{\textit{#1}}\else\rmfamily\textit{#1}\fi}}

\makeatletter

\def\extractmathfunctionname#1{\extractmathfunctionname@#1(,)\tmpa\tmpb}
\def\extractmathfunctionname@#1(#2)#3\tmpb{\def\mathname{#1}}

\makeatother
  
\newenvironment{math-function}[1]{
  \def\mathdefaultname{#1}
  \extractmathfunctionname{#1}
  \edef\mathurl{{math:\mathname}}\expandafter\hypertarget\expandafter{\mathurl}{}%
  \begin{pgfmanualentry}
    \pgfmanualentryheadline{\texttt{#1}}%
    \index{\mathname @\protect\texttt{\mathname} math function}%
    \index{Math functions!\mathname @\protect\texttt{\mathname}}%
    \pgfmanualbody
}
{
  \end{pgfmanualentry}
}

\def\pgfmanualemptytext{}
\def\pgfmanualvbarvbar{\char`\|\char`\|}

\newenvironment{math-operator}[4][]{%
  \begin{pgfmanualentry}
  \csname math#3operator\endcsname{#2}{#4}
  \def\mathtest{#4}%
  \ifx\mathtest\pgfmanualemptytext%
    \def\mathtype{(#3 operator)}
  \else%
    \def\mathtype{(#3 operator; uses the \texttt{#4} function)}
  \fi%
  \pgfmanualentryheadline{\mathexample\hfill\mathtype}%
  \def\mathtest{#1}%
  \ifx\mathtest\pgfmanualemptytext%
    \index{#2@\protect\texttt{#2} #3 math operator}%  
    \index{Math operators!#2@\protect\texttt{#2}}%
  \fi%
  \pgfmanualbody
}
{\end{pgfmanualentry}}

\newenvironment{math-operators}[5][]{%
  \begin{pgfmanualentry}
  \csname math#4operator\endcsname{#2}{#3}
  \def\mathtest{#5}%
  \ifx\mathtest\pgfmanualemptytext%
    \def\mathtype{(#4 operators)}
  \else%
    \def\mathtype{(#4 operators; use the \texttt{#5} function)}
  \fi%
  \pgfmanualentryheadline{\mathexample\hfill\mathtype}%
  \def\mathtest{#1}%
  \ifx\mathtest\pgfmanualemptytext%
    \index{#2#3@\protect\texttt{#2\protect\ #3} #4 math operators}% 
    \index{Math operators!#2#3@\protect\texttt{#2\protect\ #3}}%
  \fi%
  \pgfmanualbody
}
{\end{pgfmanualentry}}

\def\mathinfixoperator#1#2{%
  \def\mathoperator{\texttt{#1}}%
  \def\mathexample{\mvar{x}\space\texttt{#1}\space\mvar{y}}%
}

\def\mathprefixoperator#1#2{%
  \def\mathoperator{\texttt{#1}}%
  \def\mathexample{\texttt{#1}\mvar{x}}%
}

\def\mathpostfixoperator#1#2{%
  \def\mathoperator{\texttt{#1}}
  \def\mathexample{\mvar{x}\texttt{#1}}%
}

\def\mathgroupoperator#1#2{%
  \def\mathoperator{\texttt{#1\ #2}}%
  \def\mathexample{\texttt{#1}\mvar{x}\texttt{#2}}%
}

\expandafter\let\csname matharray accessoperator\endcsname=\mathgroupoperator
\expandafter\let\csname matharrayoperator\endcsname=\mathgroupoperator

\def\mathconditionaloperator#1#2{%
  \def\mathoperator{#1\space#2}
  \def\mathexample{\mvar{x}\ \texttt{#1}\ \mvar{y}\ {\texttt{#2}}\ \mvar{z}}
}

\newcommand\mathcommand[1][\mathdefaultname]{%
  \expandafter\makemathcommand#1(\empty)\stop%
  \expandafter\extractcommand\mathcommandname\@@%
  \medskip
}
\makeatletter

\def\makemathcommand#1(#2)#3\stop{%
  \expandafter\def\expandafter\mathcommandname\expandafter{\csname pgfmath#1\endcsname}%
  \ifx#2\empty%
  \else%
    \@makemathcommand#2,\stop,
  \fi}
\def\@makemathcommand#1,{%
  \ifx#1\stop%
  \else%
    \expandafter\def\expandafter\mathcommandname\expandafter{\mathcommandname{\ttfamily\char`\{#1\char`\}}}%
    \expandafter\@makemathcommand%
  \fi}
\makeatother

\def\calcname{\textsc{calc}}
%% MW: END MATH MACROS


\def\extractcommand#1#2\@@{%
  \pgfmanualentryheadline{\declare{\texttt{\string#1}}#2}%
  \removeats{#1}%
  \index{\strippedat @\protect\myprintocmmand{\strippedat}}
  \pgfmanualpdflabel{\textbackslash\strippedat}{}%
}


% \begin{environment}{{name}\marg{arguments}}
\newenvironment{environment}[1]{
  \begin{pgfmanualentry}
    \extractenvironement#1\@@
    \pgfmanualbody
}
{
  \end{pgfmanualentry}
}

\def\extractenvironement#1#2\@@{%
  \pgfmanualentryheadline{{\ttfamily\char`\\begin\char`\{\declare{#1}\char`\}}#2}%
  \pgfmanualentryheadline{{\ttfamily\ \ }\meta{environment contents}}%
  \pgfmanualentryheadline{{\ttfamily\char`\\end\char`\{\declare{#1}\char`\}}}%
  \index{#1@\protect\texttt{#1} environment}%
  \index{Environments!#1@\protect\texttt{#1}}
  \pgfmanualpdflabel{#1}{}%
}


\newenvironment{plainenvironment}[1]{
  \begin{pgfmanualentry}
    \extractplainenvironement#1\@@
    \pgfmanualbody
}
{
  \end{pgfmanualentry}
}

\def\extractplainenvironement#1#2\@@{%
  \pgfmanualentryheadline{{\ttfamily\declare{\char`\\#1}}#2}%
  \pgfmanualentryheadline{{\ttfamily\ \ }\meta{environment contents}}%
  \pgfmanualentryheadline{{\ttfamily\declare{\char`\\end#1}}}%
  \index{#1@\protect\texttt{#1} environment}%
  \index{Environments!#1@\protect\texttt{#1}}%
}


\newenvironment{contextenvironment}[1]{
  \begin{pgfmanualentry}
    \extractcontextenvironement#1\@@
    \pgfmanualbody
}
{
  \end{pgfmanualentry}
}

\def\extractcontextenvironement#1#2\@@{%
  \pgfmanualentryheadline{{\ttfamily\declare{\char`\\start#1}}#2}%
  \pgfmanualentryheadline{{\ttfamily\ \ }\meta{environment contents}}%
  \pgfmanualentryheadline{{\ttfamily\declare{\char`\\stop#1}}}%
  \index{#1@\protect\texttt{#1} environment}%
  \index{Environments!#1@\protect\texttt{#1}}}


\newenvironment{shape}[1]{
  \begin{pgfmanualentry}
    \pgfmanualentryheadline{Shape {\ttfamily\declare{#1}}}%
    \index{#1@\protect\texttt{#1} shape}%
    \index{Shapes!#1@\protect\texttt{#1}}
    \pgfmanualpdflabel{#1}{}%
    \pgfmanualbody
}
{
  \end{pgfmanualentry}
}

\newenvironment{shading}[1]{
  \begin{pgfmanualentry}
    \pgfmanualentryheadline{Shading {\ttfamily\declare{#1}}}%
    \index{#1@\protect\texttt{#1} shading}%
    \index{Shadings!#1@\protect\texttt{#1}}
  \pgfmanualpdflabel{#1}{}%
    \pgfmanualbody
}
{
  \end{pgfmanualentry}
}


\newenvironment{dataformat}[1]{
  \begin{pgfmanualentry}
    \pgfmanualentryheadline{Format {\ttfamily\declare{#1}}}%
    \index{#1@\protect\texttt{#1} format}%
    \index{Formats!#1@\protect\texttt{#1}}
  \pgfmanualpdflabel{#1}{}%
    \pgfmanualbody
}
{
  \end{pgfmanualentry}
}

\newenvironment{handler}[1]{
  \begin{pgfmanualentry}
    \extracthandler#1\@nil%
    \pgfmanualbody
}
{
  \end{pgfmanualentry}
}

\def\gobble#1{}
\def\extracthandler#1#2\@nil{%
  \pgfmanualentryheadline{Key handler \meta{key}{\ttfamily/\declare{#1}}#2}%
  \index{\gobble#1@\protect\texttt{#1} handler}%
  \index{Key handlers!#1@\protect\texttt{#1}}
  \pgfmanualpdflabel{/handlers/#1}%
}


\makeatletter


\newenvironment{stylekey}[1]{
  \begin{pgfmanualentry}
    \def\extrakeytext{style, }
    \extractkey#1\@nil%
    \pgfmanualbody
}
{
  \end{pgfmanualentry}
}

\def\choicesep{$\vert$}%
\def\choicearg#1{\texttt{#1}}

\newif\iffirstchoice

% \mchoice{choice1,choice2,choice3}
\newcommand\mchoice[1]{%
  \begingroup
  \firstchoicetrue
  \foreach \mchoice@ in {#1} {%
    \iffirstchoice
      \global\firstchoicefalse
    \else
      \choicesep
    \fi
    \choicearg{\mchoice@}%
  }%
  \endgroup
}%

% \begin{key}{/path/x=value}
% \begin{key}{/path/x=value (initially XXX)}
% \begin{key}{/path/x=value (default XXX)}
\newenvironment{key}[1]{
  \begin{pgfmanualentry}
    \def\extrakeytext{}
    %\def\altpath{\emph{\color{gray}or}}%
    \extractkey#1\@nil%
    \pgfmanualbody
}
{
  \end{pgfmanualentry}
}

% \insertpathifneeded{a key}{/pgf} -> assign mykey={/pgf/a key}
% \insertpathifneeded{/tikz/a key}{/pgf} -> assign mykey={/tikz/a key}
%
% #1: the key
% #2: a default path (or empty)
\def\insertpathifneeded#1#2{%
  \def\insertpathifneeded@@{#2}%
  \ifx\insertpathifneeded@@\empty
    \def\mykey{#1}%
  \else
    \insertpathifneeded@#2\@nil
    \ifpgfutil@in@
      \def\mykey{#2/#1}%
    \else
      \def\mykey{#1}%
    \fi
  \fi
}%
\def\insertpathifneeded@#1#2\@nil{%
  \def\insertpathifneeded@@{#1}%
  \def\insertpathifneeded@@@{/}%
  \ifx\insertpathifneeded@@\insertpathifneeded@@@
    \pgfutil@in@true
  \else
    \pgfutil@in@false
  \fi
}%

% \begin{keylist}[default path]
%   {/path/option 1=value,/path/option 2=value2}
% \end{keylist}
\newenvironment{keylist}[2][]{%
  \begin{pgfmanualentry}
    \def\extrakeytext{}%
  \foreach \xx in {#2} {%
    \expandafter\insertpathifneeded\expandafter{\xx}{#1}%
    \expandafter\extractkey\mykey\@nil%
  }%
  \pgfmanualbody
}{%
  \end{pgfmanualentry}
}%

\def\extractkey#1\@nil{%
  \pgfutil@in@={#1}%
  \ifpgfutil@in@%
    \extractkeyequal#1\@nil
  \else%
    \pgfutil@in@{(initial}{#1}%
    \ifpgfutil@in@%
      \extractequalinitial#1\@nil%
    \else
      \pgfmanualentryheadline{{\ttfamily\declare{#1}}\hfill(\extrakeytext no value)}%
      \def\mykey{#1}%
      \def\mypath{}%
      \def\myname{}%
      \firsttimetrue%
      \decompose#1/\nil%
    \fi
  \fi%
}

\def\extractkeyequal#1=#2\@nil{%
  \pgfutil@in@{(default}{#2}%
  \ifpgfutil@in@%
    \extractdefault{#1}#2\@nil%
  \else%
    \pgfutil@in@{(initial}{#2}%
    \ifpgfutil@in@%
      \extractinitial{#1}#2\@nil%
    \else
      \pgfmanualentryheadline{{\ttfamily\declare{#1}=}#2\hfill(\extrakeytext no default)}%
    \fi%
  \fi%
  \def\mykey{#1}%
  \def\mypath{}%
  \def\myname{}%
  \firsttimetrue%
  \decompose#1/\nil%
}

\def\extractdefault#1#2(default #3)\@nil{%
  \pgfmanualentryheadline{{\ttfamily\declare{#1}\opt{=}}\opt{#2}\hfill (\extrakeytext default {\ttfamily#3})}%
}

\def\extractinitial#1#2(initially #3)\@nil{%
  \pgfmanualentryheadline{{\ttfamily\declare{#1}=}#2\hfill (\extrakeytext no default, initially {\ttfamily#3})}%
}

\def\extractequalinitial#1 (initially #2)\@nil{%
  \pgfmanualentryheadline{{\ttfamily\declare{#1}}\hfill (\extrakeytext initially {\ttfamily#2})}%
  \def\mykey{#1}%
  \def\mypath{}%
  \def\myname{}%
  \firsttimetrue%
  \decompose#1/\nil%
}

% Introduces a key alias '/#1/<name of current key>'
% to be used inside of \begin{key} ... \end{key}
\def\keyalias#1{\vspace{-3pt}\item{\small alias {\ttfamily/#1/\myname}}\vspace{-2pt}\par
  \pgfmanualpdflabel{/#1/\myname}{}%
}

\newif\iffirsttime

\makeatother

\def\decompose/#1/#2\nil{%
  \def\test{#2}%
  \ifx\test\empty%
    % aha.
    \index{#1@\protect\texttt{#1} key}%
    \index{\mypath#1@\protect\texttt{#1}}%
    \def\myname{#1}%
  \else%
    \iffirsttime
      \begingroup 
        % also make a pdf link anchor with full key path.
        \def\hyperlabelwithoutslash##1/\nil{%
          \pgfmanualpdflabel{##1}{}%
        }%
        \hyperlabelwithoutslash/#1/#2\nil
      \endgroup
      \def\mypath{#1@\protect\texttt{/#1/}!}%
      \firsttimefalse
    \else
      \expandafter\def\expandafter\mypath\expandafter{\mypath#1@\protect\texttt{#1/}!}%
    \fi
    \def\firsttime{}
    \decompose/#2\nil%
  \fi%
}

\def\indexkey#1{%
  \def\mypath{}%
  \decompose#1/\nil%
}

\newenvironment{predefinedmethod}[1]{
  \begin{pgfmanualentry}
    \extractpredefinedmethod#1\@nil
    \pgfmanualbody
}
{
  \end{pgfmanualentry}
}
\def\extractpredefinedmethod#1(#2)\@nil{%
  \pgfmanualentryheadline{Method \declare{\ttfamily #1}\texttt(#2\texttt) \hfill(predefined for all classes)}
  \index{#1@\protect\texttt{#1} method}%
  \index{Methods!#1@\protect\texttt{#1}}
  \pgfmanualpdflabel{#1}{}%
}


\newenvironment{ooclass}[1]{
  \begin{pgfmanualentry}
    \def\currentclass{#1}
    \pgfmanualentryheadline{Class \declare{\texttt{#1}}}
    \index{#1@\protect\texttt{#1} class}%
    \index{Class #1@Class \protect\texttt{#1}}%
    \index{Classes!#1@\protect\texttt{#1}}
    \pgfmanualpdflabel{#1}{}%
    \pgfmanualbody
}
{
  \end{pgfmanualentry}
}

\newenvironment{method}[1]{
  \begin{pgfmanualentry}
    \extractmethod#1\@nil
    \pgfmanualbody
}
{
  \end{pgfmanualentry}
}
\def\extractmethod#1(#2)\@nil{%
  \def\test{#1}
  \ifx\test\currentclass
    \pgfmanualentryheadline{Constructor \declare{\ttfamily #1}\texttt(#2\texttt)}
  \else
    \pgfmanualentryheadline{Method \declare{\ttfamily #1}\texttt(#2\texttt)}
  \fi
  \index{#1@\protect\texttt{#1} method}%
  \index{Methods!#1@\protect\texttt{#1}}
  \index{Class \currentclass!#1@\protect\texttt{#1}}%
  \pgfmanualpdflabel{#1}{}%
}

\newenvironment{attribute}[1]{
  \begin{pgfmanualentry}
    \extractattribute#1\@nil
    \pgfmanualbody
}
{
  \end{pgfmanualentry}
}
\def\extractattribute#1=#2;\@nil{%
  \def\test{#2}%
  \ifx\test\@empty
    \pgfmanualentryheadline{Private attribute \declare{\ttfamily #1} \hfill (initially empty)}
  \else
    \pgfmanualentryheadline{Private attribute \declare{\ttfamily #1} \hfill (initially {\ttfamily #2})}
  \fi
  \index{#1@\protect\texttt{#1} attribute}%
  \index{Attributes!#1@\protect\texttt{#1}}
  \index{Class \currentclass!#1@\protect\texttt{#1}}%
  \pgfmanualpdflabel{#1}{}%
}



\newenvironment{predefinednode}[1]{
  \begin{pgfmanualentry}
    \pgfmanualentryheadline{Predefined node {\ttfamily\declare{#1}}}%
    \index{#1@\protect\texttt{#1} node}%
    \index{Predefined node!#1@\protect\texttt{#1}}
    \pgfmanualpdflabel{#1}{}%
    \pgfmanualbody
}
{
  \end{pgfmanualentry}
}

\newenvironment{coordinatesystem}[1]{
  \begin{pgfmanualentry}
    \pgfmanualentryheadline{Coordinate system {\ttfamily\declare{#1}}}%
    \index{#1@\protect\texttt{#1} coordinate system}%
    \index{Coordinate systems!#1@\protect\texttt{#1}}
    \pgfmanualpdflabel{#1}{}%
    \pgfmanualbody
}
{
  \end{pgfmanualentry}
}

\newenvironment{snake}[1]{
  \begin{pgfmanualentry}
    \pgfmanualentryheadline{Snake {\ttfamily\declare{#1}}}%
    \index{#1@\protect\texttt{#1} snake}%
    \index{Snakes!#1@\protect\texttt{#1}}
    \pgfmanualbody
}
{
  \end{pgfmanualentry}
}

\newenvironment{decoration}[1]{
  \begin{pgfmanualentry}
    \pgfmanualentryheadline{Decoration {\ttfamily\declare{#1}}}%
    \index{#1@\protect\texttt{#1} decoration}%
    \index{Decorations!#1@\protect\texttt{#1}}
    \pgfmanualbody
}
{
  \end{pgfmanualentry}
}


\def\pgfmanualbar{\char`\|}
\makeatletter
\newenvironment{pathoperation}[3][]{
  \begin{pgfmanualentry}
    \pgfmanualentryheadline{\textcolor{gray}{{\ttfamily\char`\\path}\
        \ \dots}
      \declare{\texttt{#2}}#3\ \textcolor{gray}{\dots\texttt{;}}}%
    \def\pgfmanualtest{#1}%
    \ifx\pgfmanualtest\@empty%
      \index{#2@\protect\texttt{#2} path operation}%
      \index{Path operations!#2@\protect\texttt{#2}}%
      \pgfmanualpdflabel{#2}{}%
    \fi%
    \pgfmanualbody
}
{
  \end{pgfmanualentry}
}
\newenvironment{datavisualizationoperation}[3][]{
  \begin{pgfmanualentry}
    \pgfmanualentryheadline{\textcolor{gray}{{\ttfamily\char`\\datavisualization}\
        \ \dots}
      \declare{\texttt{#2}}#3\ \textcolor{gray}{\dots\texttt{;}}}%
    \def\pgfmanualtest{#1}%
    \ifx\pgfmanualtest\@empty%
      \index{#2@\protect\texttt{#2} (data visualization)}%
      \index{Data visualization!#2@\protect\texttt{#2}}%
      \pgfmanualpdflabel{#2}{}%
    \fi%
    \pgfmanualbody
}
{
  \end{pgfmanualentry}
}
\makeatother

\def\extractcommand#1#2\@@{%
  \pgfmanualentryheadline{\declare{\texttt{\string#1}}#2}%
  \removeats{#1}%
  \index{\strippedat @\protect\myprintocmmand{\strippedat}}
  \pgfmanualpdflabel{\textbackslash\strippedat}{}%
}


\newenvironment{filedescription}[1]{
  \begin{pgfmanualentry}
    \pgfmanualentryheadline{File {\ttfamily\declare{#1}}}%
    \index{#1@\protect\texttt{#1} file}%
    \index{Packages and files!#1@\protect\texttt{#1}}%
    \pgfmanualbody
}
{
  \end{pgfmanualentry}
}


\newenvironment{packageoption}[1]{
  \begin{pgfmanualentry}
    \pgfmanualentryheadline{{\ttfamily\char`\\usepackage[\declare{#1}]\char`\{pgf\char`\}}}
    \index{#1@\protect\texttt{#1} package option}%
    \index{Package options for \textsc{pgf}!#1@\protect\texttt{#1}}%
    \pgfmanualbody
}
{
  \end{pgfmanualentry}
}



\newcommand\opt[1]{{\color{black!50!green}#1}}
\newcommand\ooarg[1]{{\ttfamily[}\meta{#1}{\ttfamily]}}

\def\pfgopt{\afterassignment\pgfmanualopt\let\next=}


{
  \makeatletter
  \global\let\myempty=\@empty
  \global\let\mygobble=\@gobble
  \catcode`\@=12
  \gdef\getridofats#1@#2\relax{%
    \def\getridtest{#2}%
    \ifx\getridtest\myempty%
      \expandafter\def\expandafter\strippedat\expandafter{\strippedat#1}
    \else%
      \expandafter\def\expandafter\strippedat\expandafter{\strippedat#1\protect\printanat}
      \getridofats#2\relax%
    \fi%
  }

  \gdef\removeats#1{%
    \let\strippedat\myempty%
    \edef\strippedtext{\stripcommand#1}%
    \expandafter\getridofats\strippedtext @\relax%
  }
  
  \gdef\stripcommand#1{\expandafter\mygobble\string#1}
}

\def\printanat{\char`\@}

\def\declare{\afterassignment\pgfmanualdeclare\let\next=}
\def\pgfmanualdeclare{\ifx\next\bgroup\bgroup\color{red!75!black}\else{\color{red!75!black}\next}\fi}


\let\textoken=\command
\let\endtextoken=\endcommand

\def\myprintocmmand#1{\texttt{\char`\\#1}}

\def\example{\par\smallskip\noindent\textit{Example: }}
\def\themeauthor{\par\smallskip\noindent\textit{Theme author: }}


\def\indexoption#1{%
  \index{#1@\protect\texttt{#1} option}%
  \index{Graphic options and styles!#1@\protect\texttt{#1}}%
}

\def\itemcalendaroption#1{\item \declare{\texttt{#1}}%
  \index{#1@\protect\texttt{#1} date test}%
  \index{Date tests!#1@\protect\texttt{#1}}%
}

\makeatletter

\ifx\scantokens\@undefined
  \PackageError{pgfmanual-macros}{You need to use extended latex
    (elatex) or (pdfelatex) to process this document}{}
\fi

\begingroup
\catcode`|=0
\catcode`[= 1
\catcode`]=2
\catcode`\{=12
\catcode `\}=12
\catcode`\\=12 |gdef|find@example#1\end{codeexample}[|endofcodeexample[#1]]
|endgroup

% define \returntospace.
%
% It should define NEWLINE as {}, spaces and tabs as \space.
\begingroup
\catcode`\^=7
\catcode`\^^M=13
\catcode`\^^I=13
\catcode`\ =13%
\gdef\returntospace{\catcode`\ =13\def {\space}\catcode`\^^I=13\def^^I{\space}}
\gdef\showreturn{\show^^M}
\endgroup

\begingroup
\catcode`\%=13
\catcode`\^^M=13
\gdef\commenthandler{\catcode`\%=13\def%{\@gobble@till@return}}
\gdef\@gobble@till@return#1^^M{}
\gdef\@gobble@till@return@ignore#1^^M{\ignorespaces}
\gdef\typesetcomment{\catcode`\%=13\def%{\@typeset@till@return}}
\gdef\@typeset@till@return#1^^M{{\def%{\char`\%}\textsl{\char`\%#1}}\par}
\endgroup

% Define tab-implementation functions
%   \codeexample@tabinit@replacementchars@
% and
%   \codeexample@tabinit@catcode@
%
% They should ONLY be used in case that tab replacement is active.
%
% This here is merely a preparation step.
%
% Idea:
% \codeexample@tabinit@catcode@ will make TAB active
% and
% \codeexample@tabinit@replacementchars@ will insert as many spaces as
% /codeexample/tabsize contains.
{
\catcode`\^^I=13
% ATTENTION: do NOT use tabs in these definitions!!
\gdef\codeexample@tabinit@replacementchars@{%
 \begingroup
 \count0=\pgfkeysvalueof{/codeexample/tabsize}\relax
 \toks0={}%
 \loop
 \ifnum\count0>0
  \advance\count0 by-1
  \toks0=\expandafter{\the\toks0\ }%
 \repeat
 \xdef\codeexample@tabinit@replacementchars@@{\the\toks0}%
 \endgroup
 \let^^I=\codeexample@tabinit@replacementchars@@
}%
\gdef\codeexample@tabinit@catcode@{\catcode`\^^I=13}%
}%

% Called after any options have been set. It assigns
%   \codeexample@tabinit@catcode
% and
%   \codeexample@tabinit@replacementchars
% which are used inside of 
%\begin{codeexample}
% ...
%\end{codeexample}
%
% \codeexample@tabinit@catcode  is either \relax or it makes tab
% active.
%
% \codeexample@tabinit@replacementchars is either \relax or it inserts
% a proper replacement sequence for tabs (as many spaces as
% configured)
\def\codeexample@tabinit{%
  \ifnum\pgfkeysvalueof{/codeexample/tabsize}=0\relax
    \let\codeexample@tabinit@replacementchars=\relax
    \let\codeexample@tabinit@catcode=\relax
  \else
    \let\codeexample@tabinit@catcode=\codeexample@tabinit@catcode@
    \let\codeexample@tabinit@replacementchars=\codeexample@tabinit@replacementchars@
  \fi
}

\pgfqkeys{/codeexample}{%
  width/.code=  {\setlength\codeexamplewidth{#1}},
  graphic/.code=  {\colorlet{graphicbackground}{#1}},
  code/.code= {\colorlet{codebackground}{#1}},
  execute code/.is if=code@execute,
  code only/.code=  {\code@executefalse},
  pre/.store in=\code@pre,
  post/.store in=\code@post,
  % #1 is the *complete* environment contents as it shall be
  % typeset. In particular, the catcodes are NOT the normal ones.
  typeset listing/.code=  {#1},
  render instead/.store in=\code@render,
  vbox/.code= {\def\code@pre{\vbox\bgroup\setlength{\hsize}{\linewidth-6pt}}\def\code@post{\egroup}},
  ignorespaces/.code= {\let\@gobble@till@return=\@gobble@till@return@ignore},
  leave comments/.code= {\def\code@catcode@hook{\catcode`\%=12}\let\commenthandler=\relax\let\typesetcomment=\relax},
  tabsize/.initial=0,% FIXME : this here is merely used for indentation. It is just a TAB REPLACEMENT.
  every codeexample/.style={width=4cm+7pt},
}

\let\code@pre\pgfutil@empty
\let\code@post\pgfutil@empty
\let\code@render\pgfutil@empty
\def\code@catcode@hook{}

\newdimen\codeexamplewidth
\newif\ifcode@execute
\newbox\codeexamplebox
\def\codeexample[#1]{%
  \begingroup%
  \code@executetrue
  \pgfqkeys{/codeexample}{every codeexample,#1}%
  \codeexample@tabinit% assigns \codeexample@tabinit@[catcode,replacementchars]
  \parindent0pt
  \begingroup%
  \par%
  \medskip%
  \let\do\@makeother%
  \dospecials%
  \obeylines%
  \@vobeyspaces%
  \catcode`\%=13%
  \catcode`\^^M=13%
  \code@catcode@hook%
  \codeexample@tabinit@catcode
  \relax%
  \find@example}
\def\endofcodeexample#1{%
  \endgroup%
  \ifcode@execute%
    \setbox\codeexamplebox=\hbox{%
      \ifx\code@render\pgfutil@empty%
      {%
        {%
          \returntospace%
          \commenthandler%
          \xdef\code@temp{#1}% removes returns and comments
        }%
        \catcode`\^^M=9%
        \colorbox{graphicbackground}{\color{black}\ignorespaces%
          \code@pre\expandafter\scantokens\expandafter{\code@temp\ignorespaces}\code@post\ignorespaces}%
      }%
      \else%
        \colorbox{graphicbackground}{\color{black}\ignorespaces%
          \code@render}%
      \fi%
    }%
    \ifdim\wd\codeexamplebox>\codeexamplewidth%
      \def\code@start{\par}%
      \def\code@flushstart{}\def\code@flushend{}%
      \def\code@mid{\parskip2pt\par\noindent}%
      \def\code@width{\linewidth-6pt}%
      \def\code@end{}%
    \else%
      \def\code@start{%
        \linewidth=\textwidth%
        \parshape \@ne 0pt \linewidth
        \leavevmode%
        \hbox\bgroup}%
      \def\code@flushstart{\hfill}%
      \def\code@flushend{\hbox{}}%
      \def\code@mid{\hskip6pt}%
      \def\code@width{\linewidth-12pt-\codeexamplewidth}%
      \def\code@end{\egroup}%
    \fi%
    \code@start%
    \noindent%
    \begin{minipage}[t]{\codeexamplewidth}\raggedright
      \hrule width0pt%
      \footnotesize\vskip-1em%
      \code@flushstart\box\codeexamplebox\code@flushend%
      \vskip-1ex
      \leavevmode%
    \end{minipage}%
  \else%
    \def\code@mid{\par}
    \def\code@width{\linewidth-6pt}
    \def\code@end{}
  \fi%
  \code@mid%  
  \colorbox{codebackground}{%
    \begin{minipage}[t]{\code@width}%
      {%
        \let\do\@makeother
        \dospecials
        \frenchspacing\@vobeyspaces
        \normalfont\ttfamily\footnotesize
        \typesetcomment%
        \codeexample@tabinit@replacementchars
        \@tempswafalse
        \def\par{%
          \if@tempswa
          \leavevmode \null \@@par\penalty\interlinepenalty
          \else
          \@tempswatrue
          \ifhmode\@@par\penalty\interlinepenalty\fi
          \fi}%
        \obeylines
        \everypar \expandafter{\the\everypar \unpenalty}%
        \pgfkeysvalueof{/codeexample/typeset listing/.@cmd}{#1}\pgfeov
      }
    \end{minipage}}%
  \code@end%
  \par%
  \medskip
  \end{codeexample}
}

\def\endcodeexample{\endgroup}


%%% Local Variables: 
%%% mode: latex
%%% TeX-master: "beameruserguide"
%%% End: 

% \hypersetup{hidelinks}
\hypersetup{%
  colorlinks=true, % use true to enable colors below:
  linkcolor=blue,%red,
  filecolor=blue,%magenta,
  urlcolor=blue,%cyan,
  citecolor=blue,
  pdfborder=0 0 0,
}

\makeindex
\makeatletter
\def\convertto#1#2{\strip@pt\dimexpr #2*65536/\number\dimexpr 1#1}%\dimexpr require etex
\input{pgfmanual.code} 
\pgfkeys{
  % list is tried to find a match.
  /pdflinks/search key prefixes in={/msc/,/tikz/,/pgf/,/depgraph/},
  %
}
\newcommand{\msckey}[3]{
 \IfEndWith*{\mscget{#1}}{pt}{%
   \IfStrEq*{#1}{label distance}{
     \def\tmp{\convertto{ex}{\mscget{#1}}ex}
    }{
     \def\tmp{\convertto{cm}{\mscget{#1}}cm}
     }
   }{%
     \edef\tmp{\mscget{#1}}
  }
  
 \begin{key}{/msc/#1=#2 (initially \tmp)}
     #3
 \end{key}
}
\newcommand{\mscstyle}[2]{
 \begin{stylekey}{/msc/#1}
     #2
 \end{stylekey}
}

% The following code is taken from the doc package. It defines a global
% macro \bslash that produces a bslash (if present in the current font).
{\catcode`\|=\z@ \catcode`\\=12 |gdef|bslash{\}}
\makeatother
\title{
  A \LaTeX{} macro package for Message Sequence Charts\\{\large User Manual}
}

\author{
 \begin{tabular}{c}
  \begin{tabular}{ccc}
   Victor Bos & &
   Ton van Deursen\\
   & & \\
   Piotr Kordy& &
   Sjouke Mauw
  \end{tabular}\\
 \end{tabular}
}

\date{\small Version A\mscversion, last update \today\\
      Describing \mscpack{} version \mscversion}

\begin{document}

\maketitle

\begin{abstract}
\noindent
The \mscpack{} facilitates the \LaTeX\ user to easily include
Message Sequence Charts in his texts. This document describes the
design and use of the \mscpack.
\end{abstract}

\tableofcontents
\section{New}
\label{new}
\paragraph{Version~2.00} is the result of a complete rewrite of the package.
The main difference is that the package is
now implemented in \tikzname\ instead of \verb+PSTricks+, making it
compatible with
\verb+pdfLaTex+ and \beamer.

In addition the syntax has been slightly extended. It now supports the
specification of all options as a comma separated list in square
brackets, while remaining backward compatible. All lengths are now
stored using \pgfname\ keys. As a result it is possible to change
sizes of shapes individually. Support for colors is added.

We are much indebted to Adrian Farmadin (Masaryk University, Czech
Republic) for initiating and contributing to this software rehaul.


\paragraph{Version~1.17} slightly modifies the handling of the
|offset| argument of the |\mess|, |\setstoptimer|, and
|\settimeout| commands. All commands with an |offset|
argument can now properly handle real numbers as
argument. Furthermore, the \verb+\topheaddist+,
\verb+\firstlevelheight+, and \verb+\instheadheight+ lengths are now
taken into account when computing the height of the instance and the
\verb+\gatesymbolradius+ length is now properly implemented.

The |\condition*| macro is now extended to dynamically
resize a condition symbol that spans multiple instances. The
|\condition*| and |\action*| macros, introduced in the
previous release, no longer shrink to a size smaller than their
non-starred counterparts. The |\condition| and
|\action| have crossed-out counterparts |\ncondition| and
|\naction| respectively.

A variation to the \verb+\stop+ command is implemented. Rather than
drawing a cross to denote that an instance has stopped, the
\verb+\stop*+ command draws an instance footer.

Finally, the \verb+\create+ command no longer requires the
\verb+creator+ parameter to be specified. If it is not specified, no
arrow is drawn from the creator to the dynamically created instance.


\paragraph{Version~1.16} solves a bug that was due to a change in
syntax of the \verb+scalebox+ macro in the  \verb+PSTricks+ package.

The |\action| and |\condition| macros are extended with a
starred option. The starred version of the macros automatically
adjusts the size of the rectangle and hexagon based on the size of the
contents.

\paragraph{Version~1.13} has a reimplementation of message commands in MSC
diagrams. The affected commands are: |\create|, |\found|,
|\lost|, |\mess|, and |\order|. The new implementation
provides more control over the placement of message labels.

The command \verb+\selfmesslabelpos+ has been removed.

The bounding box bug has been partly solved. Now, a white |\fbox|
is drawn around every (msc, mscdoc, hmsc) diagram. This makes it
possible for \texttt{dvips -E} to compute the correct bounding box for
the diagrams. Due to the |\fbox|, each diagram is extended with
0.3pt on each side (left, top, right, and bottom). This bugfix fails
if the background is not white. The xdvi program shows the white
|\fbox| in black with the result that diagrams have two visible
frames. This seems to be a bug of xdvi.

The lines around comments (|\msccomment|, Section~\ref{comments})
are changed from gray into black. The reason for this is that the gray
lines became invisible after converting the document to HTML.

\section{Introduction}
\label{introduction}

The \MSC{} language is a visual language for the description of the
interaction between different components of a system.
This language is standardized by the ITU (International
Telecommunication Union) in Recommendation Z.120~\cite{z120}. An
introductory text on \MSC{} can be found
in~\cite{RudolphGrabowskiGraubmann96}. \MSC{}s have a
wide application domain, ranging from requirements specification to
testing and documentation.
An example of a Message Sequence Chart is in Figure~\ref{quick}.

In order to support easy drawing of \MSC{}s in \LaTeX\ documents, we have
developed the \mscpack.  The current version of the \mscpack{}
supports the following \MSC{} constructs: \MSC{} frame, instances (both
single line width and double line width), messages (including
self-messages and messages to the environment), actions, singular and
combined timer events (set, timeout, reset, set-timeout, set-stop),
lost and found messages, generalized order, conditions, coregions,
activation regions, suspension regions, gates, instance creation,
instance stop, time measurements, references, and inline
expressions. In addition, there is support for \HMSC{}'s (high-level
\MSC{}s) and \MSCdoc{}s.

In this manual we explain the design and the use of the \mscpack. For
a complete overview of all features, we refer to the reference
manual~\cite{BM02}, which is included in the distribution under the
name \verb+refman.ps+. Another way to learn how to use the \mscpack{}
is to have a look at the \LaTeX{} source code of the manual and the
source code of the reference manual. They are included in the
distribution under the names \verb+manual.tex+ and \verb+refman.tex+,
respectively.  The \MSC{} constructs are simply introduced as
syntactic constructs. This paper is not meant to describe their use or
meaning.

We list the backgrounds of the package and some design decisions in
Section~\ref{background}.
Section~\ref{install} contains notes on installing the
package. Section~\ref{quickstart} contains an example of using the
package. It allows the impatient reader to quickly start using the
package. The details of using the package are explained in
Section~\ref{use}. In Section~\ref{parameters} the parameters are
explained which determine detailed layout of the various symbols.
A large but meaningless example is given in
Section~\ref{example}.


\section{Background and motivation}
\label{background}
Several commercial and non-commercial tools are available, which
support drawing or generating Message Sequence Charts. However, these
tools are in general not freely available and often not flexible
enough to satisfy all user's wishes with respect to the layout and
graphical appearance of an \MSC{}.
Therefore, people often use general drawing tools, such as {\em xfig}
to draw \MSC{}s. However flexible this approach is, it takes quite some
effort to produce nice \MSC{} drawings in a tool which is not dedicated
to \MSC{}s. Furthermore, when drawing a number of \MSC{}s it requires some
preciseness in order to obtain a consistent set of \MSC{}s.

For these reasons, we have started the design of a set of
\LaTeX\ macros which support the drawing of \MSC{}s. In this way, an \MSC{} can be
represented in \LaTeX\ in a textual format and compiled into
e.g.\ PostScript.

We aimed at satisfying the following requirements and design
decisions.
\begin{itemize}
\item
The package should follow the ITU standard with respect to shape and
placement of the symbols. (The current version supports
the \MSC2000 standard.)
\item
Static and dynamic semantics are not considered. The user is allowed
to violate all semantical restrictions and draw inconsistent \MSC{}s.
The package only supports elementary syntactical requirements.
\item
The package should offer functionality at the right level of
abstraction. Rather than supplying coordinates of pixels, the
user should be able to express the placement of symbols in terms of
{\em levels}.
Nevertheless, the textual representation of \MSC{}s as defined by the ITU
standard has a level of abstraction which is too high for our purposes.
It lacks information about the actual positioning of the \MSC{} symbols,
while we think that in our package this should be under user control.
\item
There should be only minimal automatic restructuring and layout of the
\MSC{} (e.g.\ the relative positioning of two messages should be as
defined by the user, even if the messages are not causally ordered).
\item
The user can customize the appearance of the \MSC{}s by manipulating
an appropriate set of parameters.
\end{itemize}

\section{Installation, copyright and system requirements}
\label{install}

The \mscpack{} is still under development. The authors
appreciate any comments and suggestions for improvements. The most
recent version of the package can be downloaded from
\url{http://satoss.uni.lu/mscpackage/}.


The \mscpack{} has a \emph{LaTeX Project Public License}
(LPPL), see~\url{http://www.latex-project.org/lppl.txt}:
{\small
\verbatiminput{COPYRIGHT}
} As such, it is free of charge and can be freely
distributed. Furthermore, it is allowed to make modifications to the
package, provided that modified versions get different names.  The
authors accept no liability with respect to the functioning of the
package.

The \mscpack{} runs with \LaTeXe. It has been tested with \LaTeXe\
version dated 1998/06/01 using \TeX\ version 3.14159.  The following
additional packages are required: \textsf{tikz}, \textsf{calc},
and \textsf{xstring}.  These packages are in general
part of the standard \LaTeXe\ distribution. These additional packages
can be obtained from the {\em ctan} database for \LaTeX\, e.g.\ via
the following URL: \url{http://www.tex.ac.uk/}. The \textsf{tikz}
package is described shortly in~\cite{mertz2007graphics} or in the \pgfname
manual~\cite{tantau2022pgf}.
The generated output can be previewed with pdf/ps previewing software.
It may be needed to update all \LaTeX\ related software
to a more recent version in order to smoothly work with the \mscpack.

The \mscpack{} can be installed easily. Just put the file
|msc.sty| in a directory which is searched by \LaTeX{} for style
files. The set of directories actually searched depends on the \TeX\
installation, but often the {\em current directory} is included. UNIX
users may have to set the environment variable \texttt{\$TEXINPUTS} to
an appropriate value.  For more details on this topic consult
documentation of your \TeX\ installation.

%The following text is currently not applicable -sm
%
%A problem may occur when using more packages than the \MSC{} macro
%package in your \LaTeX\ document. This happens, e.g.\ when using the
%{\em epsfig} package. These problems are due to the combination of
%several packages which are required (but not provided) by the \MSC{}
%macro package. The solution is to change the order of the
%\verb+usepackage+ clauses.

\section{Quick start}
\label{quickstart}
The \mscpack{} is easy to use.
Below is an example of the use of the package and
Figure~\ref{quick} shows the generated \MSC{}.
\begin{figure}[htbp]
\begin{center}
\begin{codeexample}[width=\linewidth-6.8cm]
\begin{msc}{Example}
\declinst{usr}{User}{}
\declinst{m1}{Machine 1}{control}
\declinst{m2}{Machine 2}{drill}

\mess[label position=below]
{startm1}{usr}{m1}
\nextlevel
\mess{startm2}{m1}{m2}
\nextlevel
\mess{log}{m1}{envleft}
\nextlevel
\nextlevel
\mess{output}{m2}{usr}[2]
\mess{free}{m1}{usr}
\nextlevel
\end{msc}
\end{codeexample}
\end{center}
\caption{The generated \MSC{}}
\label{quick}
\end{figure}
The \mscpack{} is activated by the clause \verb+\usepackage{msc}+. This package
contains, among others, a definition of the environment |msc|. This
environment is used to draw \MSC{}s. The \MSC{} definition is surrounded by the
clauses |\begin{msc}{Example}| and |\end{msc}|. The name of the
\MSC{}, {\tt Example}, is displayed in the upper-left corner of the \MSC{}.

The next four lines define the \emph{instances}:
|\declinst{m1}{Machine 1}{control}| defines an instance with
\meta{nickname} |m1| and a description consisting of two parts,
namely, |Machine 1| and |control|.  The nickname is used in
all subsequent references to this instance.  The first part of the
description is drawn above the rectangular instance head symbol, and
the second part of the description is drawn inside the instance head
symbol.

The following lines contain the definitions of the messages. Every
message has a source and a destination instance. The clause
|\mess[label position=below]{startm1}{usr}{m1}| defines a message with name
|startm1|, which goes from instance |usr| to instance |m1|.
Between square brackets we can put optional list of options. Each option is
comma separated from another option. In this example we specify the placement of
the label to be below the message arrow.

In order to control the vertical placement of the messages, the \MSC{} is
divided into levels. At every level, any number of messages may start.
The vertical position of the end point of a message is determined
by the optional fourth argument of the message definition, as in the
clause |\mess{output}{m2}{user}[2]|.
This argument is the vertical offset (in number of levels) between the
start point of the message (i.e.,\ the current level) and its end point.
If the value is 0 the message is drawn horizontally. A negative offset
means that the arrow has an upward slope. Equivalently we could specify this
value using option:  |\mess[level shift=2]{output}{m2}{user}|.

The clause |\nextlevel| is used to advance to the next level.

\section{Use of the \mscpack}
\label{use}
\subsection{Options}
Each command supported by the \mscpack\ can be provided with a list of
options. A list of
options is surrounded by square brackets and options are separated
by a comma. In fact options are \pgfname\ keys and the parsing of options
is done internally by the \tikzname\ package. All keys defined by the
\mscpack\ are
prefixed by the |/msc| keyword. When a key is not fully specified (i.e.,\ it is
specified without the |/msc|, |/tikz| or some other prefix) the
package will search first for the keys in the |/msc| namespace and if no
key is found it will search the default |/tikz| name space.

In the description below, optional arguments are marked in \opt{green}.
\subsection{The \MSC{} frame}
\label{mscframe}
The |msc| environment is used for making \MSC{} definitions. Thus, such a
definition looks as follows.
\begin{environment}{{msc}\opt{\oarg{options}}\marg{msc name}}
This declares the |msc| environment. The frame and the header of the
\MSC{} are drawn.
The argument \meta{msc name} is the name of the \MSC{}.
The header of an \MSC{} is formed from the keyword |msc|, followed
by the \meta{msc name}.

The size of the \MSC{} frame is determined vertically by the number of levels
occurring in the \MSC{} (see Section~\ref{levels}) and horizontally by
the number of instances (see Section~\ref{instances}).
\msckey{title position}{\meta{position}}{
  Controls the position of the header of the MSC. Possible values of
  \meta{position} are:
  |left|, |right|, and |center|.}
\msckey{msc keyword}{\meta{keyword}}{
  Specifies the \meta{keyword} in the header.}
\msckey{title distance}{\meta{distance}}{
  Specifies the \meta{distance} between left/right side of the frame and the MSC
  header.}
\msckey{title top distance}{\meta{distance}}{
  Specifies the \meta{distance} between top of the frame and the (top
  of the) MSC header.}
\msckey{left environment distance}{\meta{distance}}{
  Specifies the \meta{distance} between the left-most instance and the
  left side of the frame.
}
\msckey{right environment distance}{\meta{distance}}{
  Specifies the \meta{distance} between the right-most instance and
  the right side of the frame.
}
 \begin{key}{/msc/environment distance}
  Sets the keys |right environment distance| and |left environment distance| to
  the provided value.
 \end{key}
\msckey{instance distance}{\meta{distance}}{
  Sets the \meta{distance} between the (vertical) instance lines.
}

\begin{key}{/msc/draw frame=\meta{colour} (default , initially \mscget{draw frame})}
  Sets the colour of the frame. If you do not want to draw a frame, set
  \meta{colour} to |none|. When \meta{colour} is empty, then the last colour is used,
  usually black.
\end{key}

\begin{key}{/msc/draw grid=\meta{grid type} (default grid, initially \mscget{draw grid})}
  Specifies if we want to draw the help grid in msc. Possible values for
  \meta{grid type} are |grid|,
  |color grid| and |none|.
\end{key}

\end{environment}

\subsection{Levels}
\label{levels}
An \MSC{} is vertically divided in {\em levels}. All events in an \MSC{} are
attached to a certain level, or stretch out over several levels.
Any number of events can be drawn at a certain level.
An event will always be drawn (or started) at the current level, unless a level
offset is specified (see e.g.\ the |\mess| command in
Section~\ref{messages}). The \meta{level offset} is a positive integer number, which
denotes at which level, relative to the current level, an event should
be drawn.
Drawing starts at level 0.
The following command is used to advance to the next level.
\begin{command}{\nextlevel\opt{\oarg{level offset}}}
The \meta{lever offset} is a positive integer value which denotes the number of
levels to advance. By default, the value of \meta{lever offset} is 1, which means
drawing continues at the next level.
\end{command}
The following keys influence the |\nextlevel| command:
\msckey{first level height}{\meta{distance}}{
  Specifies the \meta{distance} between the instance start symbol and the first level.}
\msckey{level height}{\meta{distance}}{
  Specifies the \meta{distance} between two consecutive levels.}
\msckey{last level height}{\meta{distance}}{
  Specifies the \meta{distance} between the last level and the instance end symbol.}

Figure~\ref{parametersfig} on page~\pageref{parametersfig} shows all lengths of
the \mscpack.
\subsection{Instances}
\label{instances}
All instances have to be declared before they can be used. An instance
consists of an instance head symbol with an associated name, an
instance axis and an instance end symbol. Normal instances have a
single line axis. Fat instances have a double line axis.  The order of
the instance declarations determines the order in which the instances
occur in the drawing.

An instance is declared with the following command.

\begin{command}{\declinst\opt{*}\marg{nickname}\marg{instance name above}\marg{instance name within}}
The starred version produces a fat instance.
The \meta{nickname} is used for referring to this instance in the rest of the
MSC definition.
The \meta{instance name above} is put above the instance head symbol.
The \meta{instance name within} is put inside the instance head symbol.
\MSC{} definition.
\end{command}
The keys associated with the command:
\msckey{head height}{\meta{height}} {
  Specifes the minimum \meta{height} of the instance head symbol.}
\msckey{instance width}{\meta{width}}{
  Specifes the minimum \meta{width} of the instance head symbol.}
\msckey{head top distance}{\meta{distance}}{
  Specifes the \meta{distance} between top of the head symbols and frame.}
\msckey{foot height}{\meta{height}}{
  Specifes the \meta{height} of the instance foot symbol.}
\msckey{foot distance}{\meta{distance}}{
  Specifes the \meta{distance} between the instance foot symbol and the bottom of the \MSC{} frame.}
\msckey{environment distance}{\meta{distance}}{
  Specifies the \meta{distance} between the edge of the \MSC{} and the first/last instance axis.}
\msckey{instance width}{\meta{width}}{
  Specifies the \meta{width} of the instance head and foot symbols.}
\msckey{label distance}{\meta{distance}}{
  Specifies the \meta{distance} between the instance head symbol and the part of the
  instance name drawn above the head symbol.}
\msckey{draw head}{\meta{colour}}{
  Sets the colour of the instance head. If you do not want to draw an instance head, set
  \meta{colour} to |none|. When \meta{colour} is empty, then the last colour is used,
  usually black.
}

% The command \verb+\setfootcolor{color}+ sets the
% fill color of the footer symbol. Valid color values are, e.g.,
% \texttt{black} (default), \texttt{white}, and \texttt{gray}.
The following \MSC{} shows the declaration of an \MSC{} with three
instances. The first and the last are normal instances (one line axis),
whereas the second is a fat instance (double line axis). The optional
parameter |values small|, indicates that the small drawing
style should be used (see Section~\ref{parameters}).
% \setmscvalues{small}
\begin{codeexample}[width=\linewidth-8.8cm]
\begin{msc}[small values]{instances}
\declinst{i}{above}{within}
\declinst*{j}{}{j}
\declinst{k}{k}{}
\end{msc}
\end{codeexample}

\subsection{Messages}
\label{messages}
A message is denoted by an arrow from the sending instance to the
receiving instance. The instances are referred to by their nicknames.
A message is defined with the following command.
\begin{command}{\mess\opt{*}\opt{\oarg{options}}\marg{name}\marg{sender}\marg{receiver}\opt{\oarg{leveloffset}}}
% \begin{verbatim}
% \mess(*)[pos]{name}[labelpos]{sender}[placement]{receiver}[leveloffset]
% \end{verbatim}
The \meta{name} of the message may be any string. The \mscpack{}
processes the \meta{name} argument in LR-mode,
see~\cite[page~36]{Lam94}. This means that the string will consist of
one line. To generate multi-line message names, use the standard
\cmd{parbox} command, see the \emph{Tricks} section in the reference
manual~\cite{BM02}.  By default, the name of a message label is drawn
above the center of the arrow, but the optional keys |label position|,
|side|, and |pos| influence the actual location,
as described below.  The arrow starts at the current level at the
sending instance. The arrow ends at the current level plus the
level \meta{leveloffset}, at the receiving instance. \meta{leveloffset} can be
also modified by specifying keys |level shift| and |offset|. The
\meta{sender} and \meta{receiver} should be the nicknames of declared
instances.
Messages to or from the environment (i.e.,\ the left or the right side
of the \MSC{} frame) can be
specified by setting the sender or the receiver argument to one of the values
|envleft| or |envright|. (Note: Since instances and
environments are treated equally in the implementation, at every position
where the nickname of an instance is required, also
|envleft| and |envright| are allowed.)

In case the sending and the receiving instance are the same, the
message is a \emph{self message}. A self message is drawn as a
polyline connecting the instance axis to itself.

The starred version of the command, |\mess*|, produces the same
result as |\mess|, except that the arrow is drawn with a dashed
line. This can be used to draw a {\em method reply} (see~\cite{z120}).
\end{command}
The following keys influence the messages:
\msckey{level shift}{\meta{integer}}{
  For non-self message, \meta{integer} determines at which level message should arrive relative to the starting level
  of the message. 
}
\msckey{offset}{\meta{integer}}{
  For self message, \meta{integer} determines how many levels the message arrow should advance.
}
\msckey{side}{\meta{side}}{
  Determine the position of the arrow with respect to the instance axis. Valid values
  for \meta{side} are |left| and |right|. In case of a non-self message the key is ignored.
}
\msckey{label position}{\meta{position}}{
 Determines the position of the label.In case of non-self message, valid values
 for \meta{position} are |above|, |below|, |above~left|, |above~right|,
 |below~left| and |below~right|. In case of self message valid values are |left|
 and |right| and the default value is taken from the value of the key |side|. As
 the implementation uses corresponding |/tikz| keys the |label~position=| part
 can be skipped.
}
\msckey{pos}{\meta{real}}{
  Determines the relative distance of the message label to the beginning of the
  message. Valid values for \meta{real} are real numbers in the closed interval $[0,1]$.
  The key works in similar way as |/tikz/pos| key.
}
\msckey{self message width}{\meta{width}}{
  Controls the \meta{width} of the polyline used for drawing self-message.
%   It also affects other commands such as |\order|, |\settimeout| or |\setstoptimer|.
}
\msckey{label distance}{\meta{distance}}{
  Specifies the \meta{distance} between the label of a message and the message arrow.
}
\msckey{arrow scale}{\meta{scale}}{
  The \meta{scale} specifies by which factor should the message arrow be scaled. Influences all
  arrows drawn in |msc| environment.
}

In Figure~\ref{fig:mscsomeExample} we see \MSC{} that shows an example of the use of messages.
In this sample \MSC{} and the following \MSC{}s in this section we will
not list the complete textual representations of the \MSC{}s. For brevity
we omit the environment call and the declarations of the instances.
Note the final |\nextlevel| command which is needed
to make the instance axis long enough to receive message
\verb+a+.

Figure~\ref{fig:refpoints:B} shows different options that can be used to
position label. We use the keys |pos|, |label position| and |side| to place
the labels. Additionally we use |/tikz/anchor| key, which is an alternative way
to specify the placement of the label (see PGF Manual~\cite{tantau2022pgf}).

\begin{figure}[htbp]
\begin{codeexample}[width=\linewidth-9.8cm]
\begin{msc}[small values,head top distance=0.8cm]{messages}
\declinst{i}{}{i}
\declinst{j}{}{j}
\declinst{k}{}{k}
\mess{a}{j}{i}[3]
\mess{self}{i}{i}[1]
\nextlevel
\mess*{b}{j}{k}
\mess[b]{c}{k}{envright}
\nextlevel
\mess{d}{k}[.6]{i}
\end{msc}
\end{codeexample}
\caption{Example showing msc messages}
\label{fig:mscsomeExample}
\end{figure}
%
\begin{figure}[htbp]
\begin{codeexample}[width=\linewidth-9.8cm]
\begin{msc}[small values]{Label reference points}
\declinst{m0}{I0}{}
\declinst{m1}{I1}{}
\declinst{m2}{I2}{}
\nextlevel

\mess[label position=above]{above}{m0}{m1}
\nextlevel
\mess[label position=below]{below}{m1}{m2}
\nextlevel[2]

\mess[anchor=south,pos=0.1]{south}{m1}{m0}
\nextlevel
\mess[anchor=north]{north}{m2}{m1}
\nextlevel[2]

\mess[label position=left]{left}{m0}{m0}[2]
\mess[side=right,pos=0.9]{right}{m2}{m2}[2]
\nextlevel[4]
\mess[label position=right]{right}{m0}{m0}[2]
\mess[left,side=right]{left}{m2}{m2}[2]
\nextlevel[6]

\mess[anchor=east]{east}{m0}{m0}[-2]
\mess[anchor=west,side=right]{west}{m2}{m2}[-2]
\nextlevel[4]
\mess[right]{right}{m0}{m0}[-2]
\mess[side=right,left]{left}{m2}{m2}[-2]
\nextlevel[2]

\mess[below left]{below left}{m0}{m1}[2]
\mess[above right]{above right}{m1}{m2}[2]
\nextlevel[6]

\mess[anchor=south west]{SW}{m1}{m0}[-2]
\mess[anchor=north east,pos=0]{NE}{m2}{m1}[-2]
\nextlevel[2]
\end{msc}
\end{codeexample}
\caption{Different placing the message labels}
\label{fig:refpoints:B}
\end{figure}

\subsection{Comments}
\label{comments}

Comments are additional texts to clarify (events on) an instance. The
following command can be used to add comments to an \MSC{} diagram.
\begin{command}{\msccomment\opt{\oarg{options}}\marg{text}\marg{instance}}

The \meta{instance} parameter defines the instance to which the
comment is attached. The text of the comment is specified by the
\meta{text} parameter and is processed in LR-mode.
The key |side| affect the horizontal position relative to its instance.
\end{command}
The following key affects the drawing of the comment:
\msckey{msccomment distance}{\meta{distance}}{
  Specifies the \meta{distance} at which the comment will be placed from the
  \meta{instance} instance.
}

The following diagram shows how to use comments. In this diagram we modify the
distance between the frame and the instances in order to fit the comments inside
the frame.
\begin{codeexample}[width=\linewidth-8.8cm]
\begin{msc}[small values,
 left environment distance=2.6cm,
 right environment distance=1.9cm]{comments}
\declinst{i}{}{i}
\declinst{j}{}{j}
\mess{a}{i}{j}[2]
\msccomment[msccomment distance=1.2cm]{start}{i}
\nextlevel[2]
\msccomment[side=right]{end}{j}
\nextlevel
\end{msc}
\end{codeexample}



\subsection{Actions}
\label{actions}
An instance can perform an action, which is denoted by a
rectangle. There are two commands to draw actions:
\begin{command}{\action\opt{*}\opt{\oarg{options}}\marg{name}\marg{instance}}
The action is attached at the current level to the \meta{instance}.
The \meta{name} is centered inside the action symbol and is processed
in LR-mode. The height and width of the action symbol are adjusted so that
the \meta{name} fits inside a rectangle.

The starred version of the command, |\action*|, produces the same
result as |\action|, except that the width of the action is not adjusted.

The following keys determine the detailed drawing of the action
symbol:
\msckey{action width}{\meta{width}}{Specifies the \meta{width} of the action symbol.}
\msckey{action height}{\meta{hegiht}}{Specifies the \meta{height} of the action symbol.}
\end{command}
\begin{command}{\naction\opt{*}\opt{\oarg{options}}\marg{name}\marg{instance}}
The |\naction| command produces an action box with a diagonal
cross in it. It is otherwise identical to the |\action| command.
\end{command}


The following example shows that after an action often a multiple level
increment is needed to obtain nice pictures.
\begin{codeexample}[width=\linewidth-8.8cm]
\begin{msc}[small values, instance distance=1.6cm]{action}
\declinst{i}{}{i}
\declinst{j}{}{j}
\mess{a}{j}{i}
\nextlevel
\naction{doit}{i}
\nextlevel[2]
\action*{doitnow}{i}
\nextlevel[2]
\mess{b}{i}{j}
\end{msc}
\end{codeexample}

\subsection{Timers}
\label{timers}
There are five commands to draw timer events.
\begin{command}{\settimer\opt{\oarg{options}}\marg{name}\marg{instance}}%
Setting of a timer is drawn as a line connecting the \meta{instance}
to the {\em hour glass} symbol. The \meta{name} is put near this
symbol. The key |side| is determines on which side of the instance axis the
timer is drawn. The key |self message width| determines the lenght of the arm
between the symbol and the instance axis. The key |label distance| determines
the distance between the name and the timer symbol. The key |arrow scale|
determines the scaling of the arrow.
{\em hour glass} symbol to the \meta{instance}.
\end{command}
\begin{command}{\timeout\opt{\oarg{options}}\marg{name}\marg{instance}}
A time-out is represented by an arrow from an
{\em hour glass} symbol to the \meta{instance}.
\end{command}
\begin{command}{\stoptimer\opt{\oarg{options}}\marg{name}\marg{instance}}
Stopping a timer is drawn as a line connecting the \meta{instance} with the timer stop
symbol (denoted by a cross).
\end{command}
\begin{command}{\settimeout\opt{\oarg{options}}\marg{name}\marg{instance}\opt{\oarg{offset}}}
The command |\settimeout| is a combination of the setting of a
timer and a time out. The \meta{offset} denotes the number of levels
between the two events. 
\end{command}
\begin{command}{\setstoptimer\opt{\oarg{options}}\marg{name}\marg{instance}\opt{\oarg{offset}}}
Likewise, |\setstoptimer| is a combination
of the setting of a timer and stopping a timer.
\end{command}
\msckey{offset}{\meta{offset}}{
  The key holds the value of \meta{offset}. It denotes the number of levels
  between the two events. 
}%
\msckey{timer width}{\meta{width}}{
  The key determines the \meta{width} of the hour glass and the time out symbols.
}

The various timer symbols are shown in the following example.

\medskip
\begin{codeexample}[width=\linewidth-9.8cm]
\begin{msc}[small values]{timers}
\declinst{i}{}{i}
\declinst{j}{}{j}
\declinst{k}{}{k}
\settimer{T,50}{j}
\setstoptimer[r]{V}{k}[6]
\nextlevel[2]
\timeout{T}{j}
\settimeout{U}{i}
\nextlevel[2]
\settimer[r]{T,20}{j}
\nextlevel[2]
\stoptimer[r]{T}{j}
\end{msc}
\end{codeexample}

\subsection{Time measurements}
\label{measurements}
There are several commands to add time measurements to an \MSC{}.
\begin{command}{
\mscmark\opt{\oarg{options}}\marg{name}\marg{instance}
}
This command attach an absolute time stamp to an event on an \meta{instance}.
The \meta{name} is the text attached to the mark symbol, which is a dashed
polyline. The position of the mark is determined by the keys |position| (above,
below) and |side| (left,right).
\end{command}
\msckey{position}{\meta{position}}{
  The key determines the position of the mark symbol relative to the marked
  event. The allowed values for \meta{position} are |above|, |mid| and |below|.
}
\begin{command}{\measure\opt{*}\opt{\oarg{options}}\marg{name}\marg{instance1}\marg{instance2}\opt{\oarg{offset}}}
  This command connects two events from \meta{instance1} and \meta{instance2}.
  The first event is at the current level. The second event is \meta{offset}
  levels lower than the first event. Not that instead giving optional argument
  \meta{offset} we can specify the key |offset|. The \meta{name} is attached to
  the measure symbol. The measure symbol can be placed ath the left or right of
  \meta{instance1}. This is determined by the value of the key |side|. In the
  starred version of the command arrow heads are at the outside of the measured
  interval.
\end{command}
In case the two events are far apart, the measure may be split in two
parts. We use two commands to draw this parts:
\begin{command}{\measurestart\opt{*}\opt{\oarg{options}}\marg{name}\marg{instance}\marg{gate}}
  Draws the first part of the measure. The points where these two parts should
  be connected are drawn by a small cicle, to which the text \meta{gate} in
  attached.
  In the starred version of the command arrow heads are at the outside of the
  measured interval.
\end{command}
\begin{command}{\measureend\opt{*}\opt{\oarg{options}}\marg{name}\marg{instance}\marg{gate}}
  Draws the first part of the measure in similar way as |\measurestart| command.
\end{command}


Several keys can be used to control the detailed layout of the time measurement
symbols. The |label distance| key specify the distance between the label of a
measurement and the measurement symbol.
\msckey{measure distance}{\meta{distance}}{
This key specifies at which \meta{distance} the measurement symbols are drawn from the
\meta{instance} line.
}
\msckey{measure symbol width}{\meta{width}}{
  Specifies the \meta{width} of the measurement triangle symbol.
}

The following example illustrates marks and measurements in an \MSC{}..
\begin{codeexample}[width=\linewidth-9.8cm]
\begin{msc}[small values]{Time measurements}
\declinst{i}{}{i}
\declinst{j}{}{j}
\declinst{k}{}{k}
\mess{m1}{i}{j}[1]
\mscmark{t=0.0}{i}
\measure{0.6}{i}{j}[4]
\nextlevel
\mscmark[position=above,side=right]{t=0.3}{j}
\nextlevel[3]
\mess{m2}{j}{k}
\measurestart*[side=right]{0.2}{k}{g}
\nextlevel[6]
\mess{m3}{k}{i}
\mscmark[position=below,side=left]{t=1.0}{i}
\measureend*[side=right]{0.2}{k}{g}
\nextlevel
\end{msc}
\end{codeexample}

\subsection{Lost and found messages}
\label{lostfound}
\begin{command}{\lost\opt{\oarg{options}}\marg{name}\marg{gate}\marg{instance}}
A lost message is denoted by an arrow starting at an instance and
ending at a filled circle.
The argument \meta{instance} determines the instance to which the
arrow is attached.  The \meta{name} of the message is put above the
message arrow. The \meta{gate} is a text associated to the circle.
The keys |side|, |label position|, |pos|, |label distance| and
|self message width|
function in similar way as in the |\mess| command
(Section~\ref{messages}). That is, |side| controls the placement of the lost or
found message with respect to the instance axis. The |label position| and |pos|
controls the placement of the \meta{name} with respect to the message arrow.
The key |self message width| determines the length of the arrow and |label
distance| specifies the distance between the name and the arrow.
\end{command}
\begin{command}{\found\opt{\oarg{options}}\marg{name}\marg{gate}\marg{instance}}
A found message is denoted by an arrow starting at an open circle and ending at
an instance. Otherwise the command is identical to the |\lost| command.
\end{command}
\msckey{lost symbol radius}{\meta{radius}}{
  Determines the \meta{radius} of the circle for lost and found messages.
}

The following example shows a found and a lost message.

\begin{codeexample}[width=\linewidth-9.8cm]
\begin{msc}[small values]{lost and found}
\declinst{i}{}{i}
\declinst{j}{}{j}
\found{m}{g}{i}
\nextlevel
\mess{p}{i}{j}
\nextlevel
\lost[side=right]{n}{}{j}
\end{msc}
\end{codeexample}


\subsection{Conditions}
\label{conditions}
There are two commands to draw condition boxes:
\begin{command}{\condition\opt{*}\opt{\oarg{options}}\marg{text}\marg{instance list}}
  A condition is denoted by a hexagon. It is used to express that the
  system has entered a certain state. A condition relates to a number of
  instances. All conditions which take part in the condition are covered
  by the condition symbol. The other instances are drawn through the
  condition symbol.
  
  The \meta{text} is placed in the center of the condition. The
  \meta{instance list} expresses which instances take part in the
  condition. It is a list of nicknames of instances, separated by
  commas. Take care not to add extra white space around the nicknames,
  since this is considered part of the nickname in \LaTeX.
  The order in which the instances are listed is immaterial.

  The starred version of the command |\condition*|, produces the same
  result as the command |\condition|, except that the width of the
  condition symbol is not adjusted to fit the contents of the hexagon.
\end{command}
\begin{command}{\ncondition\opt{*}\opt{\oarg{options}}\marg{text}\marg{instance list}}
The |\ncondition| command produces a condition with a diagonal
cross in it. It is otherwise identical to the |\condition| command.
\end{command}

There are several keys that controls the shape of the condition
symbol:
\msckey{condition height}{\meta{hegiht}}{
  Determines the \meta{height} of the condition symbol.
}
\msckey{condition overlap}{\meta{distance}}{
  The \meta{distance} determines the width of the part of the condition
  symbol which extends over the rightmost/leftmost contained instance
  axis.
}

The following example contains some conditions.
\begin{codeexample}[width=\linewidth-10.6cm]
\begin{msc}[small values]{conditions}
\declinst{i}{}{i}
\declinst{j}{}{j}
\declinst{k}{}{k}
\condition{some condition}{i,k}
\nextlevel[3]
\mess{m}{i}{j}
\action{a}{k}
\nextlevel
\condition{C}{i}
\ncondition{D}{j}
\nextlevel[2]
\condition{A, B, C}{i,j,k}
\nextlevel[2]
\end{msc}
\end{codeexample}

\subsection{Generalized ordering}
\label{ordering}
A generalized order is treated much like a regular message (see
Section~\ref{messages}). There
are three differences: a generalized order is drawn with a
dotted line, it has no label, and the arrow head is in the middle of
the line. A generalized order is defined with the following command.
\begin{command}{\order\opt{\oarg{options}}\marg{sender}\marg{receiver}\opt{\oarg{level offset}}}
The \meta{sender} and \meta{receiver} are the nicknames of the
instances which are connected by the generalized ordering symbol.
At the \meta{receiver} instance, the generalized ordering symbol ends
at the current level plus the \meta{level offset}. The
\meta{level offset} is an optional integer value, with default 0.

In case \meta{sender} and \meta{receiver} denote the same instance,
the order is a \emph{self order}.  The placement of the order arrow of
a self order is controlled by the key |side|.  It
can have values |left| (meaning that the ordering symbol is drawn
left of the instance axis) and |right| (meaning that the ordering
symbol is drawn right of the instance axis). By default it is drawn at
the left side of the instance.

The position of the arrow is controlled with the |pos| key. The key |self
message width| specify the width of the polyline used for drawing orderings on a
single instance axis. The key |arrow scale| controls the size of the arrow.

Orderings to or from the environment (i.e.,\ the left or the right side
of the \MSC{} frame) can be
specified by setting the sender or the receiver argument to the value
|envleft| or |envright|.
\end{command}

An example of a generalized order is given in the following diagram.
\begin{codeexample}[width=\linewidth-9.6cm]
\begin{msc}[head top distance=0.8cm]
           {generalized order}
\declinst{i}{}{i}
\declinst{j}{}{j}
\mess{m}{envleft}{i}
\order{i}{j}[1]
\nextlevel
\mess{k}{j}{envright}
\nextlevel
\end{msc}
\end{codeexample}

\subsection{Instance regions}
\label{regions}
A part of the instance axis can be drawn in a different style. Such a
part is called an {\em instance region}. The following
regions are supported:
{\em coregion} (the instance axis is dashed, which means that the order of the
attached instances is immaterial),
{\em suspension region} (a small rectangle with dashed left and right
sides, which denotes that the instance is suspended),
{\em activation region} (a small filled rectangle, which denotes that
the instance has control).

The following commands are used to draw an instance region.
\begin{command}{\regionstart\opt{\oarg{options}}\marg{type}\marg{instance name}}
  After this command the instance axis of \meta{intstance name} is drawn in the
  shape determined by \meta{type} starting from the current level.
  The \meta{type} can have values |coregion|, |suspension|, and |activation|.
\end{command}
\begin{command}{\regionend\marg{instance name}}
  After this command the instance axis for the \meta{instance name} is drawn in
  the usual way. In other words it ends current region.
\end{command}
The following keys control the shape of on instance region:
\msckey{region width}{
  Determines the \meta{width} of the coregion start and end symbol and the \meta{width} of the activation and suspension rectangle.
}

In the following example, several instance regions are demonstrated.
Note that the relative order of |\mess|, |\regionstart|, and
|\regionend| commands does not matter. Also notice that the
second activation region on~$i$ ends with a method reply
(which is produced by the command |\mess*|).

\medskip

\begin{codeexample}[width=\linewidth-9.6cm]
\begin{msc}[small values]{regions}
\declinst{i}{}{i}
\declinst{j}{}{j}
\declinst{k}{}{k}
\regionstart{activation}{i}   \nextlevel[3]
\regionend{i}
\mess{m}{i}{j}                \nextlevel
\regionstart{coregion}{j}     \nextlevel
\regionstart{suspension}{k}
\mess{p}{j}{k}                \nextlevel
\regionstart{activation}{i}
\mess{q}{j}{i}                \nextlevel
\regionend{j}                 \nextlevel[3]
\regionend{i}
\mess*{r}{i}{j}               \nextlevel
\mess{s}{j}{k}
\regionend{k}                 \nextlevel
\end{msc}
\end{codeexample}

\subsection{Instance creation and instance stop}
\label{createstop}
The \MSC{} language offers constructs for dynamic instance creation and
instance destruction.
An instance can dynamically create another instance by issuing a
create command. An instance creation is drawn as a dashed message
arrow. At the side of the arrow head, the instance head symbol for the
created instance is drawn.
An instance end symbol does not denote the end of the specified
process, but merely the end of its current description. Therefore, a
different symbol is needed which denotes that an instance stops before
the end of the \MSC{} in which it is contained.
The instance stop symbol is a cross.

The following commands are used for instance creation and instance
stop.
\begin{command}{\dummyinst\opt{\oarg{options}}\marg{created instance}}
  This command reserve space for on instance which will be created dynamically.
  |\dummyinst| command is mixed with the declarations of normal instances (see
  the |\declinst| command, Section~\ref{instances}). The argument \meta{created
  instance} is the nickname of the instance that will be created later.
\end{command}
\begin{command}{\startinst\opt{\oarg{options}}\marg{instance name}%
    \marg{text above}\marg{text inside}}
  This command create an instance. The created instance must have been declared
  first with the |\dummyinst| command. The name of the created instance consists
  of two parts. The part called \meta{text above} is placed above the instance
  head and the \meta{text inside} is centered within the instance head.
\end{command}
\begin{command}{\create\opt{\oarg{options}}\marg{name}\marg{creator}\marg{instance name}%
    \marg{text above}\marg{text inside}}
  This command is identical to |\startinst| command except it allows to specify
  the creator of the instance. The command results in a horizontal message arrow
  labelled with \meta{name}. The arrow starts at the current level at the
  instance with the nickname \meta{creator} and ends at the current level at the
  instance head of the instance with nickname \meta{instance name}.
  As with normal messages, placement of the message label is controlled by the
  keys |label position| and |pos|.
\end{command}
\begin{command}{\stop\opt*\opt{\oarg{options}}\marg{instance name}}
  An instance is stopped with the |\stop| command.
  The \meta{instance} is the nickname of the stopped instance. The
  instance axis is not drawn any further below the level at which the
  |\stop| command is issued. Also, the instance foot symbol is not
  drawn.
  
  If the |\stop*| command is used, then the instance foot symbol is
  drawn instead of the stop symbol.
\end{command}
\msckey{stop width}{\meta{width}}{
  The \meta{width} determines the size of the stop symbol.
}


Take care not to specify any events on an instance which has not yet
been created or which has already been stopped. However, it is possible
to create an instance after it has stopped, as showed in the next example.

\begin{codeexample}[width=\linewidth-9.6cm]
\begin{msc}[small values]{dynamic instances}
\declinst{i}{}{j}
\dummyinst{j}
\declinst{k}{}{k}

\mess{p}{i}{k}
\nextlevel[3]
\create[label position=below]{kick}{k}{j}{}{j}
\nextlevel[2]
\stop*{k}
\nextlevel
\mess{ok}{j}{i}
\nextlevel
\stop{j}
\nextlevel[2]
\startinst{k}{}{again}
\nextlevel[2]
\end{msc}
\end{codeexample}

\subsection{\MSC{} references}
\label{references}
Within an \MSC{} a reference to other \MSC{}s can be included. Such a
reference is drawn as a rectangle with rounded corners, covering part
of the \MSC{}. The following commands are used to draw \MSC{} references.

\begin{command}{\referencestart\opt{\oarg{options}}\marg{nickname}
    \marg{text}\marg{left instance}\marg{right instance}}
  The reference symbol starts at the level where the
  |\referencestart| command is used, and ends at the level where
  the corresponding |\referenceend| command occurs. These commands
  correspond if they have the same \meta{nickname}. 
  The \meta{text} is
  placed in the center of the reference symbol. The reference covers all
  instances from \meta{leftinstance} to \meta{rightinstance}.

  The left and right edge of the reference symbol are
  \meta{nickname}|left| and \meta{nickname}|right|, where \meta{nickname} is
  the nickname of the \MSC{} reference as defined in the
  |\referencestart| command.
  These names can be used at every place where the nickname of an
  instance is required, e.g.\ as the sender or receiver of a message.
\end{command}
\begin{command}{\referenceend\opt{\oarg{options}}\marg{nickname}}
  This command end the reference \meta{nickname} that has been started with the
  |\referencestart| command.
\end{command}
The following keys affect the drawing of reference symbol:
\msckey{left reference overlap}{\meta{distance}}{
  Determines the \meta{distance} between the left  edge of the reference symbol and the
  leftmost covered instance axis.
}
\msckey{right reference overlap}{\meta{distance}}{
  Determines the \meta{distance} between the right  edge of the reference symbol and the
  rightmost covered instance axis.
}

The following example shows the usage of references:

\begin{codeexample}[width=\linewidth-9.6cm]
\begin{msc}[small values]{references}
\declinst{i}{}{i}
\declinst{j}{}{j}
\declinst{k}{}{k}
\mess{a}{j}{k}
\nextlevel
\referencestart{r}{a reference}{i}{j}
\nextlevel
\mess{b}{rright}{k}
\nextlevel[2]
\referenceend{r}
\nextlevel[2]
\referencestart{p}{another reference}{i}{k}
\nextlevel
\mess{c}{envright}{pright}
\nextlevel[4]
\referenceend{p}
\end{msc}
\end{codeexample}

\subsection{Inline expressions}
\label{inlines}
An inline expression is a part of an \MSC{} on which an operation is
defined. A rectangle surrounds the part of the \MSC{} containing the
operands. The operands are separated by horizontal dashed lines.
The operator is placed in the upper left corner of the inline
expression symbol.
The following commands are used to draw inline expressions.

\begin{command}{\inlinestart\opt{\oarg{options}}\marg{nickname}\marg{operator}\marg{left instance}\marg{right instance}}
  Starts the inline expression at the current level. The inline expression
  spans over the instances from \meta{left instance} to \meta{right instance}.
  The \meta{operand} is placed in the upper left corner of the rectangle.
% 
  The left and right edge of the inline expression symbol are named
  \meta{nickname}|left| and \meta{nickname}|right|, where \meta{nickname} is
  the nickname of the inline expression as defined in the
  |\inlinestart| command.
  These names can be used at every place where the nickname of an
  instance is required, e.g.\ as the sender or receiver of a message.
\end{command}
\begin{command}{\inlineseparator\opt{\oarg{options}}\marg{nickname}}
  Draws a dashed line at the current level for the inline expression with
  nickname \meta{nickname}.
\end{command}
\begin{command}{\inlineend\opt{*}\opt{\oarg{options}}\marg{nickname}}
  Ends the inline expresion with the nickname \meta{nickname} at the current level.
  The command |\inlineend*| does the same as the command
  |\inlineend|, except that the bottom line of the rectangle is
  This is used to indicate that the operand is optional.
  dashed.
\end{command}

\msckey{left inline overlap}{\meta{distance}}{
  Determines the \meta{distance} which is a distance between the left edge of the inline expression
  symbol and the leftmost  included instance axis. 
}
\msckey{right inline overlap}{\meta{distance}}{
  Determines the \meta{distance} which is a distance between the right edge of the inline expression
  symbol and the rightmost included instance axis. 
}

The following example shows how to draw inline expressions:

\begin{codeexample}[width=\linewidth-9.6cm]
\begin{msc}[small values]{inline expressions}
\declinst{i}{}{i}
\declinst{j}{}{j}
\declinst{k}{}{k}
\inlinestart{exp1}{par}{i}{k}      \nextlevel
\mess{a}{k}{j}                     \nextlevel
\inlinestart[left inline overlap=0.6cm]
            {exp2}{alt}{i}{j}      \nextlevel[2]
\mess{b}{j}{i}                     \nextlevel
\inlineseparator{exp2}             \nextlevel
\mess{c}{j}{i}                     \nextlevel
\inlineend{exp2}                   \nextlevel
\inlineseparator{exp1}             \nextlevel
\mess{d}{j}{k}                     \nextlevel
\inlineend{exp1}                   \nextlevel
\inlinestart{exp3}{opt}{i}{j}      \nextlevel
\mess{e}{envright}{exp3right}
\mess{e}{exp3right}{j}             \nextlevel
\mess{f}{j}{i}                     \nextlevel
\inlineend*{exp3}
\end{msc}
\end{codeexample}
%
\subsection{Gates}
\label{gates}

A gate determines a connection point for messages.

The following command can be used to draw gates.

\begin{command}{\gate\opt{*}\opt{\oarg{options}}\marg{gate name}\marg{instance name}}
  The unstarred version produces a normal (invisible) gate. The starred
  version produces a visible gate (a small dot). The gate is drawn at the
  current level at the instance \meta{instance name} (which can also be the
  left and right edge of e.g.\ an \MSC{} reference).  The
  \meta{gate name} is attached to the gate. The positioning of the \meta{gate
  name} relative to the gate is determined by the keys |side| (|left| or
  |right|) and |position| (|above|, |mid| or |below|).
\end{command}

There are several keys to control the size and shape of the gate symbol:
\msckey{gate symbol radius}{\meta{radius}}{
  The \meta{radius} defines the radius of the gate symbol.
}

The next example shows a number of gates.
\begin{codeexample}[width=\linewidth-9.6cm]
\begin{msc}[small values]{gates}
\declinst{i}{}{i}
\declinst{j}{}{j}
\declinst{k}{}{k}
\referencestart{r}{ref}{i}{j}
\nextlevel
\gate{$g$}{rright}
\mess{b}{rright}{envright}
\gate[r][c]{$h$}{envright}
\nextlevel[2]
\mess{c}{rright}{k}
\gate*[l][b]{$h'$}{rright}
\nextlevel[2]
\referenceend{r}
\end{msc}
\end{codeexample}

\subsection{High-level \MSC{}s}
\label{hmsc}
A High-level \MSC{} (\HMSC{}) is a drawing which defines the relation between
a number of \MSC{}s. It is composed of a start symbol (an upside down
triangle), a number of end symbols (represented by triangles), a
number of \MSC{} references (these are rectangles with rounded corners),
a number of conditions (hexagons) and possibly several connection
points (circles). These symbols are connected by arrows.

The following commands can be used to draw \HMSC{}s.
\begin{environment}{{hmsc}\opt{\oarg{options}}\marg{hmsc name}\opt{(\meta{llx},\meta{lly})(\meta{urx},\meta{ury})}}
  In order to draw \HMSC{}s, a new environment is defined, which is called
  |hmsc|. The command to begin this environment has several
  arguments.
  
  The header of an \HMSC{} is
  formed from the value of the key |hmsc keyword| (initially |hmsc|), followed by the
  \meta{hmsc name}. The positioning of the header depends on the key
  |title position|. The possible values are |left|, |right| and |center|. The
  keys |title top distance| and |title distance| control the distances of the
  header from the frame.
  The size of the \HMSC{} frame is calculated automatically unless the optional
  parameters (\meta{llx},\meta{lly}) and (\meta{urx},\meta{ury}) are specified. In
  such case the frame is determined by the
  coordinates of the lower-left corner,(\meta{llx},\meta{lly}) , and the
  coordinates of the upper-right corner, (\meta{urx},\meta{ury}).
\msckey{hmsc keyword}{\meta{keyword}}%
{The \meta{keyword} will be a starting word in a header of \HMSC{}.}
\msckey{north hmsc margin}{\meta{distance}}%
{The \meta{distance} determines the north margin of hmsc environment}
\msckey{south hmsc margin}{\meta{distance}}%
{The \meta{distance} determines the south margin of hmsc environment}
\msckey{east hmsc margin}{\meta{distance}}%
{The \meta{distance} determines the east margin of hmsc environment}
\msckey{west hmsc margin}{\meta{distance}}%
{The \meta{distance} determines the west margin of hmsc environment}
\msckey{west hmsc margin}{\meta{distance}}%
{The \meta{distance} determines the west margin of hmsc environment}
\begin{stylekey}{/msc/hmsc margin=\meta{distance}}
     Sets the |north hmsc margin|, |south hmsc margin|, |west hmsc margin|,
     |east hmsc margin| to \meta{distance}.
\end{stylekey}
\end{environment}
In the \HMSC{} we can draw number of symbols. Each symbol is in fact a \tikz{}
node and \meta{options} can be used to specify for example the placement
relative to the other nodes. The following commands are used to draw the symbols:
\begin{command}{\hmscstartsymbol\opt{\oarg{options}}\marg{nicknamee}\opt{(\meta{x},\meta{y})}}
  Draws an equilateral triangle pointed down representing the start symbol. The
  \meta{nickname} can be used later as a reference to the drawn symbol. Optional
  \meta{x} and \meta{y} specify the coordinates at which the symbol should be
  drawn. 
\end{command}
\begin{stylekey}{/msc/start symbol}
  A style that is applied to draw a triangle node in
  the |\hmscstartsymbol| command.
\end{stylekey}
\msckey{hmsc symbol width}{\meta{width}}{
  Sets the width of the triangle of the start and end symbol in hmsc to the \meta{width}.  
}
\begin{command}{\hmscendsymbol\opt{\oarg{options}}\marg{nickname}\opt{(\meta{x},\meta{y})}}
  The command in identical to |\hmscstartsymbol| command but draws an
  equilateral triangle pointed up representing the start symbol.
\end{command}
\begin{stylekey}{/msc/end symbol}
  A style that is applied to draw a triangle node in
  the |\hmscendsymbol| command.
\end{stylekey}
\begin{command}{\hmscreference\opt{\oarg{options}}\marg{nickname}\marg{text}\opt{(\meta{x},\meta{y})}}
  Draws a \meta{text} inside a frame with rounded corners. The \meta{nickname}
  can be used later to connect it to other hmsc elements.
\end{command}
\begin{stylekey}{/msc/reference}
  A style that is applied to draw a reference node in the |\hmscreference| command.
\end{stylekey}
\msckey{reference height}{\meta{height}}%
{ Sets the minimal height of the reference to \meta{height}.}
\msckey{reference width}{\meta{width}}%
{ Sets the minimal width of the reference to \meta{width}.}
\begin{command}{\hmsccondion\opt{\oarg{options}}\marg{nickname}\marg{text}\opt{(\meta{x},\meta{y})}}
  Draws a \meta{text} inside a rectangular frame with diamond shaped sides. The \meta{nickname}
  can be used later to connect it to other hmsc elements.
\end{command}
\msckey{hmsc condition height}{\meta{height}}%
{ Sets the minimal height of the condition to \meta{height}.}
\msckey{hmsc condition width}{\meta{width}}%
{ Sets the minimal width of the condition to \meta{width}.}
\begin{command}{\hmscconnection\opt{\oarg{options}}\marg{nickname}\opt{(\meta{x},\meta{y})}}
  Draws a small circle symbolizing connection at coordinates
  (\meta{x},\meta{y}). Later we can refer to the connection using \meta{nickname}.
\end{command}
\msckey{hmsc connection radius}{\meta{distance}}%
{ Sets the radius of the connection symbol to the \meta{distance}}

The \HMSC{} grid is not drawn, but used to control the positioning of the
\HMSC{} symbols  (|startsymbol|, |endsymbol|, |reference|,
|condition|, and |connection|). 
The center of each symbol is drawn on the grid point
with coordinates (\meta{x},\meta{y}). Each symbol also has a \meta{nickname}
for later reference.

\begin{verbatim}
\hmscreference{nickname}{text}(x,y)
\hmsccondition{nickname}{text}(x,y)
\hmscconnection{nickname}(x,y)
\arrow{from-nickname}[coord-list]{to-nickname}
\end{verbatim}

\HMSC{} symbols can be connected by means of the \verb+arrow+ command.
This draws an arrow from the symbol with nickname \verb+from-nickname+
to the symbol with nickname \verb+to-nickname+. The optional argument
\verb+coord-list+ can be used if the line connecting the source and
the destination should not be straight. The \verb+coord-list+ has
the following syntax: \verb+(x1,y1)(x2,y2)...(xk,yk)+. This means
that the connecting line goes through the points with coordinates
(x1,y1), (x2,y2), \ldots, (xk,yk).

Arrows always leave the start symbol at the bottom. They enter the end symbol
at the top.
Arrows start and end either at the middle of the top or at the middle
of the bottom of a reference and condition symbol. The incoming
(outgoing) direction of the arrow determines whether it will start
(end) at the top or at the bottom.

There are several parameters to control the size and shape of the
symbols (see Section~\ref{parameters}). These are
\verb+\hmscconditionheight+ (the height of the condition symbol),
\verb+\hmscconditionwidth+ (the width of the condition symbol,
excluding the left and right angular parts),
\verb+\hmscreferenceheight+ (the height of the reference symbol),
\verb+\hmscreference width+ (the width of the reference symbol),
\verb+\messarrowscale{size}+ (a command to set the size of the arrow
head of a connection line);
\verb+setconnectiontype(type)+ (set the shape of the polyline
connection the symbols; \verb+type+ can be \verb+straight+,
\verb+rounded+, and \verb+curved+),
\verb+\startsymbolwidth+ (the width of the start and end symbol),
\verb+\topnamedist+ (sets the distance between the top of the \HMSC{} frame
and the \HMSC{} header).

An example of an \HMSC{} is in the following diagram. Notice that the
width and height of reference symbols are changed locally (i.e.,
between \verb+{+ and \verb+}+ braces) just before the big
reference~\verb+b+ is defined.

\medskip

\noindent
\begin{minipage}[c]{0.375\linewidth}

\begin{hmsc}{High level}(-3,0)(3,12)
\hmscstartsymbol{s}(0,10)
\hmscreference{a}{A}(0,9)
\hmscconnection{c}(0,7.5)
{
\mscset{reference width=2cm}
\mscset{reference height=2\baselineskip}
% \setlength{\hmscreferencewidth}{2cm}
% \setlength{\hmscreferenceheight}{3\baselineskip}
\hmscreference{b}{\parbox{1.9cm}
  {\centering
   this is a big
   hmsc reference}
  }(0,6)
}
\hmsccondition{t}{Again}(1,4.25)
\hmsccondition{ok}{Ok}(0,3.5)
\hmsccondition{q}{\textit{Break}}(-2,3.5)
\hmscendsymbol{e1}(-2,2)
\hmscreference{do}{Do it!}(0,2)
\hmscendsymbol{e2}(0,1)

\arrow{s}{a}
\arrow{a}{c}
\arrow{c}{b}
\arrow{c}[(-2,7.5)]{q}
\arrow{b}{q}
\arrow{q}{e1}
\arrow{b}{ok}
\arrow{ok}{do}
\arrow{do}{e2}
\arrow{b}{t}
\arrow{t}[(1,3.5)(2.5,3.5)(2.5,7.5)]{c}
\end{hmsc}
\end{minipage}\hfill
\begin{minipage}[c]{0.6\linewidth}
{\footnotesize
\begin{verbatim}
\begin{hmsc}{High level}(-3,0)(3,12)
\hmscstartsymbol{s}(0,10)
\hmscreference{a}{A}(0,9)
\hmscconnection{c}(0,7.5)
{\setlength{\hmscreferencewidth}{2cm}
 \setlength{\hmscreferenceheight}
           {3\baselineskip}
 \hmscreference{b}{\parbox{1.9cm}
   {\centering this is a big hmsc reference}
   }(0,6)}
\hmsccondition{t}{Again}(1,4.5)
\hmsccondition{ok}{Ok}(0,3.5)
\hmsccondition{q}{\textit{Break}}(-2,3.5)
\hmscendsymbol{e1}(-2,2)
\hmscreference{do}{Do it!}(0,2)
\hmscendsymbol{e2}(0,1)
\arrow{s}{a}
\arrow{a}{c}
\arrow{c}{b}
\arrow{c}[(-2,7.5)]{q}
\arrow{b}{q}
\arrow{q}{e1}
\arrow{b}{ok}
\arrow{ok}{do}
\arrow{do}{e2}
\arrow{b}{t}
\arrow{t}[(1,3.5)(2.5,3.5)(2.5,7.5)]{c}
\end{hmsc}
\end{verbatim}
}
\end{minipage}


\subsection{\MSC{} documents}
\label{mscdoc}
An \MSCdoc{} is a drawing which contains various declarations
of objects used in the \MSC{} description. For drawing \MSCdoc{}s
he following commands are provided.

\begin{verbatim}
\begin{mscdoc}[headerpos]{mscdocname}{text}(llx,lly)(urx,ury)
\end{mscdoc}
\reference{text}(x,y)
\separator{y}
\end{verbatim}

As for \MSC{} and \HMSC{}, a new environment is defined, which is named
\verb+mscdoc+. The command to begin an \MSCdoc{} has several
arguments. The argument \verb+headerpos+ is optional. It controls
positioning of the header of the \MSCdoc{}.
This argument can have values
\verb+l+ (for a left aligned header), \verb+c+ (for a centered header)
and \verb+r+ (for a right aligned header).
The header of an \MSCdoc{} is formed from the keyword \verb+mscdocument+,
followed by the \verb+mscdocname+.
The \verb+text+ is placed left-aligned below the \MSCdoc{} header.
The size of the \MSCdoc{} frame is determined by
coordinates of the lower-left corner, \verb+(llx,lly)+, and the
coordinates of the upper-right corner, \verb+(urx,ury)+.

The \MSCdoc{} grid is not drawn, but used to control the positioning of the
\MSC{} references.
The center of such a reference is drawn on the grid point
with coordinates \verb+{x,y}+.

The \verb+separator+ command draws a dashed horizontal line. The \MSC{}
references above the separator are the exported, while the ones below
the separator are local.

There are several parameters to control the size and shape of the
symbols (see Section~\ref{parameters}).
\verb+\mscdocreferenceheight+ (the height of the reference symbol),
\verb+\mscdocreferencewidth+ (the width of the reference symbol),
\verb+\topnamedist+ (sets the distance between the top of the \MSCdoc{}
frame and the \MSCdoc{} header).

An example of an \MSCdoc{} is in the following diagram. Notice that
the size of references in an \MSCdoc{} had to be changed for the last
reference.

\medskip


\begin{minipage}{0.4\linewidth}
\begin{mscdoc}{My declarations}%
              (0,0)(6,8)
\reference{a}(1,6)
\reference{b}(3,6)
\reference{c}(1,4)
\reference{d}(3,4)
\separator{3}
\mscset{reference width=4.5cm}
\mscset{reference height=2\baselineskip}
\reference{%
 \parbox{4cm}%
 {\raggedright This is a
  two-line reference}}
 (3,1.5)
\end{mscdoc}
\end{minipage}
%
\begin{minipage}{0.5\linewidth}
{\small
\begin{verbatim}
\begin{mscdoc}{My declarations}%
              (0,0)(6,8)
\reference{a}(1,6)
\reference{b}(3,6)
\reference{c}(1,4)
\reference{d}(3,4)
\separator{3}
\mscset{reference width=4.5cm}
\mscset{reference height=2\baselineskip}
\reference{%
 \parbox{4cm}%
 {\raggedright This is a
  two-line reference}}
 (3,1.5)
\end{mscdoc}
\end{verbatim}
}
\end{minipage}



\section{Style parameters}
\label{parameters}
\definecolor{markColor}{rgb}{0.10,0.10,0.9}

\makeatletter
\NewDocumentCommand{\measureMark}{%
    s       % arrows outside the measured distance
    D<>{1.5pt} % offset of marks
    O{.75cm}    % offset of label
    m       % first point
    m       % second point
    m       % label
    D<>{o}  % () coordinates -> angle
            % h -> horizontal,
            % v -> vertical
            % o or what ever -> oblique
    O{}     % options for tikzset
    }{%

    {\tikzset{#8}

    \coordinate (@1) at #4 ;
    \coordinate (@2) at #5 ;

    \if #7v % vertical
        \coordinate (@0) at ($($#4!.5!#5$) + (#3,0)$) ; 
        \coordinate (@4) at (@0|-@1) ;
        \coordinate (@5) at (@0|-@2) ;
    \else
    \if #7h % horizontal
        \coordinate (@0) at ($($#4!.5!#5$) + (0,#3)$) ; 
        \coordinate (@4) at (@0-|@1) ;
        \coordinate (@5) at (@0-|@2) ;
    \else % cotation encoche
    \ifnum\pdfstrcmp{\unexpanded\expandafter{\@car#7\@nil}}{(}=\z@
        \coordinate (@5) at ($#7!#3!#5$) ;
        \coordinate (@4) at ($#7!#3!#4$) ;
    \else % cotation oblique    
        \coordinate (@5) at ($#5!#3!90:#4$) ;
        \coordinate (@4) at ($#4!#3!-90:#5$) ;
    \fi\fi\fi

    \draw[dotted,very thin,shorten >= #2,shorten <= -2*#2] (@4) -- #4 ;
    \draw[dotted,very thin,shorten >= #2,shorten <= -2*#2] (@5) -- #5 ;

    \IfBooleanTF #1 {% with star
    \draw[measureMark arrow,-,draw=none] (@4) -- (@5)
        node[measureMark node] {#6\strut};
    \draw[measureMark arrow,<-] (@4) -- ($(@4)!-2pt!(@5)$) ;   
    \draw[measureMark arrow,<-] (@5) -- ($(@5)!-2pt!(@4)$) ;   
    }{% without star
    \ifnum\pdfstrcmp{\unexpanded\expandafter{\@car#7\@nil}}{(}=\z@
        \draw[measureMark arrow] (@5) to[bend right]
            node[measureMark node] {#6\strut} (@4) ;
    \else
    \draw[measureMark arrow] (@4) -- (@5)
        node[measureMark node] {#6\strut};
    \fi
    }}
    }
\makeatother

By means of a collection of parameters, the graphical appearance of
an \MSC{} can be fine tuned to the user's taste.
The general parameters
are displayed in Figure~\ref{parametersfig} on page~\pageref{parametersfig}.
%
% %========================
\begin{figure}[htb]
  \begin{msc}[right environment distance=2.3cm]{Lengths}
   \tikzset{%
       measureMark node/.style={%
           midway,
           sloped,
           fill=white,
           text = markColor,
           font=\footnotesize,
           inner sep=1.5pt,
           outer sep=2pt
       },
       measureMark arrow/.style={%
           <->,
           >=latex,
           very thin,dashed
       },
       dimen/.style={<->,>=latex,thin,every rectangle node/.style={fill=white,midway,font=\sffamily}},
       symmetry/.style={dashed,thin},
   }
  \declinst{m1}{aname}{iname}
  \declinst{st}{}{}
  \declinst{m2}{}{}
  
  \mess{msg1}{m1}{st}
  \nextlevel
  \mess{msg2}{m1}{st}
  \nextlevel
  \mess{self}{m1}{m1}
  \coregionstart{m2}
  \nextlevel[2]
  \coregionend{m2}
  \nextlevel
  \lost[r]{x}{}{m2}
  \stop{st}
  \settimer{T}{m1}
  \nextlevel
  \inlinestart[left inline overlap=0.9cm,right inline overlap=0.5cm]{alt}{alt}{m1}{m1}
  \nextlevel[2]
  \inlineend{alt}
  \action{a}{m2}
  \nextlevel
  \referencestart{ref}{\small reference}{m1}{m1}
  \nextlevel
  \referenceend{ref}
  \nextlevel
  \gate*{}{m1}
  \nextlevel
  \condition{cond1}{m1,m2}
  \nextlevel[3]
  \mess{msg3}{m2}{m1}
  \makeatletter
    \msc@endmsc[]
  \makeatother
    
    \measureMark*[-2.5cm]{($(msc@env@nw)+(\mscget{title distance},0cm)$)}{($(msc@env@nw)+
      (\mscget{title distance},-\mscget{title top distance})$)}%
    {title top distance}%
    [measureMark node/.append style={above right=-36pt and -2pt,rotate=90}]
  
    \measureMark*[-0.8cm]{($(msc@env@nw)+(0cm,-\mscget{title top distance}-1em)$)}{($(msc@env@nw)+
      (\mscget{title distance},-\mscget{title top distance}-1em)$)}%
    {title distance}[measureMark node/.append style={left=6pt}]
  
    \measureMark*[-1.1cm]{(msc@nodem1.north)}{($(msc@nodem1.north)+ (0cm,\mscget{label distance})$)}%
    {label distance}[measureMark node/.append style={above left=6pt and 9pt,rotate=-90}]
  
    \measureMark*[-1.6cm]{($(msc@level2-|msc@nodem1)+(1.5cm,0)$)}
    {($(msc@level2-|msc@nodem1)+ (1.5cm,\mscget{label distance})$)}%
    {label distance}[measureMark node/.append style={above left=6pt and 9pt,rotate=-90}]
  
    \measureMark[0.0cm]{($(msc@env@nw)+(4cm,0)$)}{(msc@nodest.north)}%
    {head top distance}<v>[measureMark node/.append style={right=0pt,rotate=90}]
  
    \measureMark[0.2cm]{($(msc@env@nw)+(0,-3cm)$)}{(msc@nodem1.south)}%
    {left environment distance}<h>[measureMark node/.append style={above left=0pt and -7pt}]
  
    \measureMark[1.0cm]{($(msc@nodem1.south)+(-\mscget{self message
        width},-2.4cm)$)}{($(msc@nodem1.south)+(0cm,-2.4cm)$)}%
    {self message width}<h>[measureMark node/.append style={left=8pt}]
  
    \measureMark*[0.21cm]{($(msc@nodem1.south)+(-\mscget{self message width}-
      \mscget{timer width}/2,-2.9cm)$)}
    {($(msc@nodem1.south)+(-\mscget{self message width}+\mscget{timer width}/2,-2.9cm)$)}%
    {timer width}<h>[measureMark node/.append style={left=7pt,yshift=1pt}]
  
    \measureMark[0.17cm]{($(msc@level8-|msc@nodem1)+(-0.9cm,-10pt)$)}
    {($(msc@level8-|msc@nodem1)+(0,-10pt)$)}%
    {left inline overlap}<h>[measureMark node/.append style={left=12pt}]
  
    \measureMark[0.17cm]{($(msc@level8-|msc@nodem1)+(0.5cm,-2pt)$)}
    {($(msc@level8-|msc@nodem1)+(0,-2pt)$)}%
    {right inline overlap}<h>[measureMark node/.append style={right=6pt}]
  
    \measureMark[0.0cm]{($(msc@level11-|msc@nodem1)+(-\mscget{left reference overlap},4pt)$)}
    {($(msc@level11-|msc@nodem1)+(0,4pt)$)}%
    {left reference overlap}<h>[measureMark node/.append style={left=14pt}]
  
    \measureMark*[-0.2cm]{($(msc@level12-|msc@nodem1)+(0,\mscget{gate symbol radius})$)}
    {($(msc@level12-|msc@nodem1)$)}%
    {gate symbol radius}<v>[measureMark node/.append style={left=1pt,rotate=90}]
    
    \measureMark[-0.3cm]{($(msc@level17-|msc@nodem1)+(-\mscget{instance width}/2,0pt)$)}
    {($(msc@level17-|msc@nodem1)+(\mscget{instance width}/2,0pt)$)}%
    {instance width}<h>[measureMark node/.append style={below=0pt}]
  
    \measureMark[0.2cm]{(msc@level0-|msc@nodest)}
    {(msc@level1-|msc@nodest)}%
    {first level height}<v>[measureMark node/.append style={align=left,right=0pt,rotate=90}]
  
    \measureMark[0.2cm]{(msc@level1-|msc@nodest)}
    {(msc@level2-|msc@nodest)}%
    {level height}<v>[measureMark node/.append style={align=left,right=0pt,rotate=90}]
  
    \measureMark[0.2cm]{(msc@nodem2.north east)}
    {(msc@nodem2.south east)}%
    {head height}<v>[measureMark node/.append style={align=left,right=0pt,rotate=90}]
  
    \measureMark[0.2cm]{($(msc@level3-|msc@nodem2)+(\mscget{region width}/2,0)$)}
    {($(msc@level3-|msc@nodem2)+(-\mscget{region width}/2,0)$)}%
    {region width}<h>[measureMark node/.append style={align=left,right=5pt}]
  
    \measureMark*[0.2cm]{($(msc@level6-|msc@nodem2)+(\mscget{self message width},\mscget{lost symbol radius})$)}
    {($(msc@level6-|msc@nodem2)+(\mscget{self message width},0pt)$)}%
    {lost symbol radius}<v>[measureMark node/.append style={align=left,right=0pt,rotate=90}]
  
    \measureMark[-0.2cm]{($(msc@level6-|msc@nodest)+(-\mscget{stop width}/2,-\mscget{stop width}/2)$)}
    {($(msc@level6-|msc@nodest)+(\mscget{stop width}/2,-\mscget{stop width}/2)$)}%
    {stop width}<h>[measureMark node/.append style={align=left,right=7pt}]
  
    \measureMark[-0.2cm]{($(msc@level4-|msc@nodem1)$)}
    {($(msc@level4-|msc@nodest)$)}%
    {instance\\ distance}<h>[measureMark node/.append style={align=left,fill=none,yshift=-1pt}]
  
    \measureMark[0.2cm]{($(msc@level9-|msc@nodem2)+(-\mscget{action width}/2,0)$)}
    {($(msc@level9-|msc@nodem2)+(\mscget{action width}/2,0)$)}%
    {action width}<h>[measureMark node/.append style={align=left,above=-1pt,xshift=3pt}]
  
    \measureMark[0.2cm]{($(msc@level9-|msc@nodem2)+(\mscget{action width}/2,0)$)}
    {($(msc@level9-|msc@nodem2)+(\mscget{action width}/2,-\mscget{action height})$)}%
    {action height}<v>[measureMark node/.append style={align=left,right=1pt,rotate=90}]
  
    \measureMark[0.5cm]{($(msc@level13-|msc@nodem2)+(\mscget{condition overlap},0)$)}
    {($(msc@level13-|msc@nodem2)+(\mscget{condition overlap},-\mscget{condition height})$)}%
    {condition height}<v>[measureMark node/.append style={align=left,right=1pt,rotate=90}]
  
    \measureMark[0.2cm]{($(msc@level13-|msc@nodem2)+(\mscget{condition overlap},0)$)}
    {($(msc@level13-|msc@nodem2)+(0,0pt)$)}%
    {right condition overlap}<h>[measureMark node/.append style={align=left,above=-1pt}]
  
    \measureMark<2pt>[0.2cm]{($(msc@level13-|msc@nodem2)+(\mscget{condition overlap},0)$)}
    {($(msc@level13-|msc@nodem2)+(0,0pt)$)}%
    {right condition overlap}<h>[measureMark node/.append style={align=left,above=-1pt}]
  
    \measureMark<2pt>[0.2cm]{($(msc@level16-|msc@nodem2)+(0,0)$)}
    {($(msc@level17-|msc@nodem2)+(0,0pt)$)}%
    {last level height}<v>[measureMark node/.append style={align=left,right=-2pt,rotate=90}]
  
    \measureMark*[0.2cm]{($(msc@level17-|msc@nodem2)+(\mscget{instance width}/2,0)$)}
    {($(msc@level17-|msc@nodem2)+(\mscget{instance width}/2,-\mscget{foot height})$)}%
    {foot height}<v>[measureMark node/.append style={align=left,right=1pt,rotate=90}]
  
    \measureMark<0pt>[0.0cm]{($(msc@level17-|msc@nodem2)+(0.3cm,-\mscget{foot height})$)}
    {($(msc@env@se-|msc@nodem2)+(0.3cm,0)$)}%
    {foot distance}<v>[measureMark node/.append style={align=left,left=-1pt,rotate=90}]
  
    \measureMark[0.2cm]{($(msc@env@se)+(0,7.5cm)$)}{(msc@nodem2.south)}%
    {right\\ environment\\ distance}<h>[measureMark node/.append style={align=left,right=33pt}]
    
    \mscset{draw frame=none,msc keyword={},msc/name={}}
  \end{msc}
  
  \caption{User controllable parameters}
  \label{parametersfig}
\end{figure}
%
% The value of a parameter can be changed using standard
% \LaTeX\ commands, e.g.
% \begin{verbatim}
% \setlength{\levelheight}{1cm}
% \end{verbatim}
%
% The following list describes all parameters.
% The default values for drawing \MSC{}s at large, normal and small size are
% included. See the command \verb+\setmscvalues{size}+ below for
% restoring the parameters to their original values.
%
% %========================
% \begin{defs}
%
% \item[\cmd{actionheight}]
% Height of action symbols.\\
% (\lnsvalue{0.75}{0.6}{0.5} cm.)
%
% \item[\cmd{actionwidth}]
% Width of action symbol.\\
% (\lnsvalue{1.25}{1.25}{1.2} cm.)
%
% \item[\cmd{bottomfootdist}]
% Distance between bottom of foot symbol and frame.\\
% (\lnsvalue{1.0}{0.7}{0.5} cm.)
%
% \item[\cmd{commentdist}]
% Distance between a comment and its instance.\\
% (\lnsvalue{0.5}{0.5}{0.5} cm.)
%
% \item[\cmd{conditionheight}]
% Height of condition symbols.\\
% (\lnsvalue{0.75}{0.6}{0.5} cm.)
%
% \item[\cmd{conditionoverlap}]
% Overlap of condition symbol.\\
% (\lnsvalue{0.6}{0.5}{0.4} cm.)
%
% \item[\cmd{envinstdist}]
% Distance between environments and nearest instance line.\\
% (\lnsvalue{2.5}{2.0}{1.2} cm.)
%
% \item[\cmd{firstlevelheight}] Height of level just below head
% symbols. Should not be changed inside the \MSC{} environment.\\
% (\lnsvalue{0.75}{0.6}{0.4} cm.)
%
% \item[\cmd{gatesymbolradius}] Radius of the gate symbol.\\
% (\lnsvalue{0.5}{0.5}{0.5} mm.)
%
% \item[\cmd{hmscconditionheight}]
% Height of \HMSC{} condition symbol.\\
% (\lnsvalue{0.375}{0.3}{0.25} cm.)
%
% \item[\cmd{hmscconditionwidth}]
% Width of \HMSC{} condition symbol.\\
% (\lnsvalue{1.0}{0.8}{0.7} cm.)
%
% \item[\cmd{hmscconnectionradius}]
% Radius of \HMSC{} connection symbol.\\
% (\lnsvalue{0.06}{0.05}{0.04} cm.)
%
% \item[\cmd{hmscreferenceheight}]
% Height of \HMSC{} reference symbol.\\
% (\lnsvalue{0.375}{0.3}{0.25} cm.)
%
% \item[\cmd{hmscreferencewidth}]
% Width of \HMSC{} reference symbol.\\
% (\lnsvalue{1.0}{0.8}{0.7} cm.)
%
% \item[\cmd{hmscstartsymbolwidth}]
% Width of \HMSC{} start symbol.\\
% (\lnsvalue{0.75}{0.6}{0.3} cm.)
%
% \item[\cmd{inlineoverlap}]
% Overlap of inline symbol.\\
% (\lnsvalue{1.5}{1.0}{0.75} cm.)
%
% \item[\cmd{instbarwidth}]
% Default width of vertical instance bars (applies to fat instances only).\\
% (\lnsvalue{0.0}{0.0}{0.0} cm.)
%
% \item[\cmd{instdist}]
% Distance between instance axes.\\
% (\lnsvalue{3.0}{2.2}{1.5} cm.)
%
% \item[\cmd{instfootheight}] Height of foot symbols. Should not be
% changed inside the \MSC{} environment.\\
% (\lnsvalue{0.25}{0.2}{0.15} cm.)
%
% \item[\cmd{instheadheight}] Height of head symbols. Should not be
% changed inside the \MSC{} environment.\\
% (\lnsvalue{0.6}{0.55}{0.5} cm.)
%
% \item[\cmd{instwidth}]
% Width of header and foot symbols.\\
% (\lnsvalue{1.75}{1.6}{1.2} cm.)
%
% \item[\cmd{labeldist}]
% Distance between labels and the symbols to which they belong (for instance, message labels and arrows).\\
% (\lnsvalue{1.0}{1.0}{1.0} ex.)
%
% \item[\cmd{lastlevelheight}] Height of level just above foot
% symbols. Should not be changed inside the \MSC{} environment.\\
% (\lnsvalue{0.5}{0.4}{0.3} cm.)
%
% \item[\cmd{leftnamedist}] Distance between left of the frame and
% (left of) \MSC, \HMSC, or \MSCdoc{} title.\\
% (\lnsvalue{0.3}{0.2}{0.1} cm.)
%
% \item[\cmd{levelheight}]
% Height of a level.\\
% (\lnsvalue{0.75}{0.5}{0.4} cm.)
%
% \item[\cmd{lostsymbolradius}]
% Radius of the lost and found symbols.\\
% (\lnsvalue{0.15}{0.12}{0.08} cm.)
%
% \item[\cmd{markdist}]
% Horizontal distance from a mark to its instance.\\
% (\lnsvalue{1.0}{1.0}{1.0} cm.)
%
% \item[\cmd{measuredist}]
% Horizontal distance from a measure to its (closest) instance.\\
% (\lnsvalue{1.0}{1.0}{1.0} cm.)
%
% \item[\cmd{measuresymbolwidth}]
% Width of a measure symbol.\\
% (\lnsvalue{0.75}{0.6}{0.4} cm.)
%
% \item[\cmd{mscdocreferenceheight}]
% Height of reference symbol in an \MSCdoc.\\
% (\lnsvalue{0.375}{0.3}{0.25} cm.)
%
% \item[\cmd{mscdocreferencewidth}]
% Width of reference symbol in an \MSCdoc.\\
% (\lnsvalue{1.0}{0.8}{0.7} cm.)
%
% \item[\cmd{referenceoverlap}]
% Overlap of reference symbol.\\
% (\lnsvalue{1.5}{1.0}{0.75} cm.)
%
% \item[\cmd{regionbarwidth}]
% Width of region bars.\\
% (\lnsvalue{0.5}{0.4}{0.2} cm.)
%
% \item[\cmd{selfmesswidth}]
% Length of horizontal arms of self messages.\\
% (\lnsvalue{0.75}{0.6}{0.4} cm.)
%
% \item[\cmd{stopwidth}]
% Width of the stop symbol.\\
% (\lnsvalue{0.6}{0.5}{0.3} cm.)
%
% \item[\cmd{timerwidth}]
% Width of the \emph{timer} symbols.\\
% (\lnsvalue{0.4}{0.3}{0.2} cm.)
%
% \item[\cmd{topheaddist}]
% Distance between top of head symbols and frame.\\
% (\lnsvalue{1.5}{1.3}{1.2} cm.)
%
% \item[\cmd{topnamedist}] Distance between top of the frame and
% (top of) \MSC, \HMSC, or \MSCdoc{} title.\\
% (\lnsvalue{0.3}{0.2}{0.2} cm.)
%
% \end{defs}
%
% In addition there are several commands which allow the user to adjust
% the \MSC{} drawing to his own taste.
%
% \begin{defs}
%
% \item[\cmd{messarrowscale}\{\emph{scalefactor}\}] Sets the scale
% factor (a positive real number) of message arrow
% heads. (\lnsvalue{2}{1.5}{1.2})
%
% \item[\cmd{setmscscale}\{\emph{scalefactor}\}] Sets the scale factor
% of the \MSC{} environment to \emph{scalefactor}. the scale factor is
% supposed to be a real number. Scaling is done when the \MSC{}
% environment ends (\verb|\end{msc}|). The default of \emph{scalefactor}
% is~1. A more consistent way for
% varying the size of the \MSC{} can be obtained by using the
% \verb+\setmesvalues+ command as described below.
% (default value 1.)
%
% \item[\cmd{psset}\{\texttt{linewidth=D}\}] This command sets the width
% of all lines in \MSC{}s, \HMSC{}s, and \MSCdoc{}s to length \verb+D+.  If
% this command is issued outside the msc environment, then the value is
% set for the complete document. If it is used directly after the start
% of the msc environment it only holds for this \MSC{}.\\
% (large/normal/small value 0.8/0.7/0.6 pt.)
%
% \item[\cmd{setfootcolor}\{\emph{color}\}] Sets the color of the foot
% symbols of \MSC{} instances. Possible values are \emph{black},
% \emph{white}, \emph{gray}, or \emph{lightgray}. For more color values,
% see the documentation of the \LaTeXe{} \textsf{color} package.
%
% \end{defs}
%
% The following command can be used to set the above mentioned
% style parameters to suitable values.
%
% \begin{defs}
% \item[\cmd{setmscvalues}\{\emph{size}\}] Sets all parameters of the \mscpack{} to
%  predefined values. Valid values for \emph{size} are:
% \verb|small|, \verb|normal|, and \verb|large|.
% (The default value of \verb+size+ is \verb+normal+.
% This can be used for drawings with at maximum six instances on a
% sheet of A4 paper.
% For sizes \verb+large+ and \verb+small+, a maximum of four and nine
% instances respectively fit on a sheet of A4 paper.)
% \end{defs}
%
% Caution has to be taken when changing the value of a parameter within
% the \MSC{} definition. The following parameters can be changed within an
% \MSC{} definition without unexpected side effects:
%
% \begin{flushleft}
% \verb+\actionheight+,
% \verb+\actionwidth+,
% \verb+\bottomfootdist+,
% \verb+\conditionheight+,
% \verb+\conditionoverlap+,
% \verb+\inlineoverlap+,
% \verb+\instfootheight+,
% \verb+\instwidth+ (however, this may cause different sizes of
% corresponding instance header and footer symbols),
% \verb+\labeldist+,
% \verb+\lastlevelheight+,
% \verb+\levelheight+,
% \verb+\lostsymbolradius+,
% \verb+\referenceoverlap+,
% \verb+\regionbarwidth+,
% \verb+\selfmesswidth+,
% \verb+\stopwidth+,
% \verb+\timerwidth+, and
% \verb+\topnamedist+.
% \end{flushleft}
%
%
% In addition to the parameters specific to the \mscpack,
% standard \LaTeX\ commands can be used to change the type style and
% other details.
% For example, if the command \verb+\sffamily+ is included directly after
% the start of the msc environment, the text in the \MSC{}
% is drawn using a {\sf sans serif} font. Likewise, the text size can be
% changed by inserting, e.g., the command \verb+\small+.
% The \verb+\raisebox+ and \verb+\parbox+ commands
% can also be used to position and format names.


\section{Example}
\label{example}
Figure~\ref{ex} on page~\pageref{ex} shows the \MSC{} defined in the
following \LaTeX\ fragment.

{\small
\begin{verbatim}
\begin{msc}{Example}
\declinst{usr}{The user}{User}
\declinst{m1}{Control}{M1}
\dummyinst{m2}
\declinst{m3}{Another Machine}{M3}

\create{start}{m1}{m2}{Processing}{M2}
\mess{msg 0}{envleft}{usr}
\mess{msg 1}{envright}{m2}[1]
\nextlevel

\mess{msg 2}{usr}{m1}
\order{m1}{m2}[4]
\action{a}{m3}

\nextlevel
\found{msg x}{}{usr}
\nextlevel

\mess{msg 3}{usr}{m2}[-1]
\coregionstart{m1}
\settimeout{S}{m3}[2]
\nextlevel


\mess{msg 4}{m1}{usr}
\coregionstart{m2}
\settimer[r]{T}{m3}
\nextlevel

\mess[r]{msg 5}{m2}{m2}[3]
\mess{msg 6}{usr}{usr}[2]
\nextlevel

\mess{msg 7}{m2}{usr}
\timeout[r]{T}{m3}
\nextlevel

\coregionend{m2}
\nextlevel

\coregionend{m1}
\stoptimer[r]{T'}{m3}
\nextlevel

\lost[r]{msg y}{Mach 1}{usr}
\mess{msg 8}{m1}{envright}
\nextlevel

\condition{condition 1}{usr,m2}
\setstoptimer[r]{U}{m3}
\nextlevel[2]
\stop{usr}

\end{msc}
\end{verbatim}
}

\begin{figure}[htb]
\begin{center}
\begin{msc}{Example}
\declinst{usr}{The user}{User}
\declinst{m1}{Control}{M1}
\dummyinst{m2}
\declinst{m3}{Another Machine}{M3}

\create{start}{m1}{m2}{Processing}{M2}
\mess{msg 0}{envleft}{usr}
\mess{msg 1}{envright}{m2}[1]
\nextlevel

\mess{msg 2}{usr}{m1}
\order{m1}{m2}[4]
\action{a}{m3}

\nextlevel
\found{msg x}{}{usr}
\nextlevel

\mess{msg 3}{usr}{m2}[-1]
\coregionstart{m1}
\settimeout{S}{m3}[2]
\nextlevel


\mess{msg 4}{m1}{usr}
\coregionstart{m2}
\settimer[r]{T}{m3}
\nextlevel

\mess[r]{msg 5}{m2}{m2}[3]
\mess{msg 6}{usr}{usr}[2]
\nextlevel

\mess{msg 7}{m2}{usr}
\timeout[r]{T}{m3}
\nextlevel

\coregionend{m2}
\nextlevel

\coregionend{m1}
\stoptimer[r]{T'}{m3}
\nextlevel

\lost[r]{msg y}{Mach 1}{usr}
\mess{msg 8}{m1}{envright}
\nextlevel

\condition{condition 1}{usr,m2}
\setstoptimer[r]{U}{m3}
\nextlevel[2]
\stop{usr}

\end{msc}

\caption{A menagerie of \MSC{} symbols}
\label{ex}
\end{center}
\end{figure}


\section{Acknowledgments}
Thanks are due to the following people for providing us with useful
input: Cas Cremers, Hugo Jonker, Peter Peters, Sa\v{s}a Radomirovi\'c,
Michel Reniers, and Pim Vullers.

\bibliographystyle{plain}
\bibliography{biblio}
\printindex

\end{document}
